\documentclass[
  doc,
  floatsintext,
  longtable,
  a4paper,
  nolmodern,
  notxfonts,
  notimes,
  12pt,
  colorlinks=true,linkcolor=blue,citecolor=blue,urlcolor=blue]{apa7}

\usepackage{amsmath}
\usepackage{amssymb}

\geometry{inner=1in, outer=1in}
\fancyhfoffset[LE,RO]{0cm}


\usepackage[bidi=default]{babel}
\babelprovide[main,import]{spanish}


% get rid of language-specific shorthands (see #6817):
\let\LanguageShortHands\languageshorthands
\def\languageshorthands#1{}

\RequirePackage{longtable}
\RequirePackage{threeparttablex}

\makeatletter
\renewcommand{\paragraph}{\@startsection{paragraph}{4}{\parindent}%
	{0\baselineskip \@plus 0.2ex \@minus 0.2ex}%
	{-.5em}%
	{\normalfont\normalsize\bfseries\typesectitle}}

\renewcommand{\subparagraph}[1]{\@startsection{subparagraph}{5}{0.5em}%
	{0\baselineskip \@plus 0.2ex \@minus 0.2ex}%
	{-\z@\relax}%
	{\normalfont\normalsize\bfseries\itshape\hspace{\parindent}{#1}\textit{\addperi}}{\relax}}
\makeatother




\usepackage{longtable, booktabs, multirow, multicol, colortbl, hhline, caption, array, float, xpatch}
\usepackage{subcaption}


\renewcommand\thesubfigure{\Alph{subfigure}}
\setcounter{topnumber}{2}
\setcounter{bottomnumber}{2}
\setcounter{totalnumber}{4}
\renewcommand{\topfraction}{0.85}
\renewcommand{\bottomfraction}{0.85}
\renewcommand{\textfraction}{0.15}
\renewcommand{\floatpagefraction}{0.7}

\usepackage{tcolorbox}
\tcbuselibrary{listings,theorems, breakable, skins}
\usepackage{fontawesome5}

\definecolor{quarto-callout-color}{HTML}{909090}
\definecolor{quarto-callout-note-color}{HTML}{0758E5}
\definecolor{quarto-callout-important-color}{HTML}{CC1914}
\definecolor{quarto-callout-warning-color}{HTML}{EB9113}
\definecolor{quarto-callout-tip-color}{HTML}{00A047}
\definecolor{quarto-callout-caution-color}{HTML}{FC5300}
\definecolor{quarto-callout-color-frame}{HTML}{ACACAC}
\definecolor{quarto-callout-note-color-frame}{HTML}{4582EC}
\definecolor{quarto-callout-important-color-frame}{HTML}{D9534F}
\definecolor{quarto-callout-warning-color-frame}{HTML}{F0AD4E}
\definecolor{quarto-callout-tip-color-frame}{HTML}{02B875}
\definecolor{quarto-callout-caution-color-frame}{HTML}{FD7E14}

%\newlength\Oldarrayrulewidth
%\newlength\Oldtabcolsep


\usepackage{hyperref}



\usepackage{color}
\usepackage{fancyvrb}
\newcommand{\VerbBar}{|}
\newcommand{\VERB}{\Verb[commandchars=\\\{\}]}
\DefineVerbatimEnvironment{Highlighting}{Verbatim}{commandchars=\\\{\}}
% Add ',fontsize=\small' for more characters per line
\usepackage{framed}
\definecolor{shadecolor}{RGB}{241,243,245}
\newenvironment{Shaded}{\begin{snugshade}}{\end{snugshade}}
\newcommand{\AlertTok}[1]{\textcolor[rgb]{0.68,0.00,0.00}{#1}}
\newcommand{\AnnotationTok}[1]{\textcolor[rgb]{0.37,0.37,0.37}{#1}}
\newcommand{\AttributeTok}[1]{\textcolor[rgb]{0.40,0.45,0.13}{#1}}
\newcommand{\BaseNTok}[1]{\textcolor[rgb]{0.68,0.00,0.00}{#1}}
\newcommand{\BuiltInTok}[1]{\textcolor[rgb]{0.00,0.23,0.31}{#1}}
\newcommand{\CharTok}[1]{\textcolor[rgb]{0.13,0.47,0.30}{#1}}
\newcommand{\CommentTok}[1]{\textcolor[rgb]{0.37,0.37,0.37}{#1}}
\newcommand{\CommentVarTok}[1]{\textcolor[rgb]{0.37,0.37,0.37}{\textit{#1}}}
\newcommand{\ConstantTok}[1]{\textcolor[rgb]{0.56,0.35,0.01}{#1}}
\newcommand{\ControlFlowTok}[1]{\textcolor[rgb]{0.00,0.23,0.31}{\textbf{#1}}}
\newcommand{\DataTypeTok}[1]{\textcolor[rgb]{0.68,0.00,0.00}{#1}}
\newcommand{\DecValTok}[1]{\textcolor[rgb]{0.68,0.00,0.00}{#1}}
\newcommand{\DocumentationTok}[1]{\textcolor[rgb]{0.37,0.37,0.37}{\textit{#1}}}
\newcommand{\ErrorTok}[1]{\textcolor[rgb]{0.68,0.00,0.00}{#1}}
\newcommand{\ExtensionTok}[1]{\textcolor[rgb]{0.00,0.23,0.31}{#1}}
\newcommand{\FloatTok}[1]{\textcolor[rgb]{0.68,0.00,0.00}{#1}}
\newcommand{\FunctionTok}[1]{\textcolor[rgb]{0.28,0.35,0.67}{#1}}
\newcommand{\ImportTok}[1]{\textcolor[rgb]{0.00,0.46,0.62}{#1}}
\newcommand{\InformationTok}[1]{\textcolor[rgb]{0.37,0.37,0.37}{#1}}
\newcommand{\KeywordTok}[1]{\textcolor[rgb]{0.00,0.23,0.31}{\textbf{#1}}}
\newcommand{\NormalTok}[1]{\textcolor[rgb]{0.00,0.23,0.31}{#1}}
\newcommand{\OperatorTok}[1]{\textcolor[rgb]{0.37,0.37,0.37}{#1}}
\newcommand{\OtherTok}[1]{\textcolor[rgb]{0.00,0.23,0.31}{#1}}
\newcommand{\PreprocessorTok}[1]{\textcolor[rgb]{0.68,0.00,0.00}{#1}}
\newcommand{\RegionMarkerTok}[1]{\textcolor[rgb]{0.00,0.23,0.31}{#1}}
\newcommand{\SpecialCharTok}[1]{\textcolor[rgb]{0.37,0.37,0.37}{#1}}
\newcommand{\SpecialStringTok}[1]{\textcolor[rgb]{0.13,0.47,0.30}{#1}}
\newcommand{\StringTok}[1]{\textcolor[rgb]{0.13,0.47,0.30}{#1}}
\newcommand{\VariableTok}[1]{\textcolor[rgb]{0.07,0.07,0.07}{#1}}
\newcommand{\VerbatimStringTok}[1]{\textcolor[rgb]{0.13,0.47,0.30}{#1}}
\newcommand{\WarningTok}[1]{\textcolor[rgb]{0.37,0.37,0.37}{\textit{#1}}}

\providecommand{\tightlist}{%
  \setlength{\itemsep}{0pt}\setlength{\parskip}{0pt}}
\usepackage{longtable,booktabs,array}
\usepackage{calc} % for calculating minipage widths
% Correct order of tables after \paragraph or \subparagraph
\usepackage{etoolbox}
\makeatletter
\patchcmd\longtable{\par}{\if@noskipsec\mbox{}\fi\par}{}{}
\makeatother
% Allow footnotes in longtable head/foot
\IfFileExists{footnotehyper.sty}{\usepackage{footnotehyper}}{\usepackage{footnote}}
\makesavenoteenv{longtable}

\usepackage{graphicx}
\makeatletter
\newsavebox\pandoc@box
\newcommand*\pandocbounded[1]{% scales image to fit in text height/width
  \sbox\pandoc@box{#1}%
  \Gscale@div\@tempa{\textheight}{\dimexpr\ht\pandoc@box+\dp\pandoc@box\relax}%
  \Gscale@div\@tempb{\linewidth}{\wd\pandoc@box}%
  \ifdim\@tempb\p@<\@tempa\p@\let\@tempa\@tempb\fi% select the smaller of both
  \ifdim\@tempa\p@<\p@\scalebox{\@tempa}{\usebox\pandoc@box}%
  \else\usebox{\pandoc@box}%
  \fi%
}
% Set default figure placement to htbp
\def\fps@figure{htbp}
\makeatother







\usepackage{newtx}

\defaultfontfeatures{Scale=MatchLowercase}
\defaultfontfeatures[\rmfamily]{Ligatures=TeX,Scale=1}





\title{Guía completa de Git y GitHub desde cero}


\shorttitle{Guía completa de Git y GitHub desde cero}


\usepackage{etoolbox}



\ccoppy{\textcopyright~2023}



\author{Edison Achalma}



\affiliation{
{Escuela Profesional de Economía, Universidad Nacional de San Cristóbal
de Huamanga}}




\leftheader{Achalma}

\date{2023-02-16}


\abstract{Este abstract será actualizado una vez que se complete el
contenido final del artículo. }

\keywords{keyword1, keyword2}

\authornote{\par{\addORCIDlink{Edison Achalma}{0000-0001-6996-3364}} 
\par{ }
\par{   El autor no tiene conflictos de interés que revelar.    Los
roles de autor se clasificaron utilizando la taxonomía de roles de
colaborador (CRediT; https://credit.niso.org/) de la siguiente
manera:  Edison Achalma:   conceptualización, redacción}
\par{La correspondencia relativa a este artículo debe dirigirse a Edison
Achalma, Email: \href{mailto:elmer.achalma.09@unsch.edu.pe}{elmer.achalma.09@unsch.edu.pe}}
}

\makeatletter
\let\endoldlt\endlongtable
\def\endlongtable{
\hline
\endoldlt
}
\makeatother

\urlstyle{same}



\makeatletter
\@ifpackageloaded{caption}{}{\usepackage{caption}}
\AtBeginDocument{%
\ifdefined\contentsname
  \renewcommand*\contentsname{Tabla de contenidos}
\else
  \newcommand\contentsname{Tabla de contenidos}
\fi
\ifdefined\listfigurename
  \renewcommand*\listfigurename{Lista de Figuras}
\else
  \newcommand\listfigurename{Lista de Figuras}
\fi
\ifdefined\listtablename
  \renewcommand*\listtablename{Lista de Tablas}
\else
  \newcommand\listtablename{Lista de Tablas}
\fi
\ifdefined\figurename
  \renewcommand*\figurename{Figura}
\else
  \newcommand\figurename{Figura}
\fi
\ifdefined\tablename
  \renewcommand*\tablename{Tabla}
\else
  \newcommand\tablename{Tabla}
\fi
}
\@ifpackageloaded{float}{}{\usepackage{float}}
\floatstyle{ruled}
\@ifundefined{c@chapter}{\newfloat{codelisting}{h}{lop}}{\newfloat{codelisting}{h}{lop}[chapter]}
\floatname{codelisting}{Listing}
\newcommand*\listoflistings{\listof{codelisting}{Lista de Listings}}
\makeatother
\makeatletter
\makeatother
\makeatletter
\@ifpackageloaded{caption}{}{\usepackage{caption}}
\@ifpackageloaded{subcaption}{}{\usepackage{subcaption}}
\makeatother
\makeatletter
\@ifpackageloaded{fontawesome5}{}{\usepackage{fontawesome5}}
\makeatother

% From https://tex.stackexchange.com/a/645996/211326
%%% apa7 doesn't want to add appendix section titles in the toc
%%% let's make it do it
\makeatletter
\xpatchcmd{\appendix}
  {\par}
  {\addcontentsline{toc}{section}{\@currentlabelname}\par}
  {}{}
\makeatother

%% Disable longtable counter
%% https://tex.stackexchange.com/a/248395/211326

\usepackage{etoolbox}

\makeatletter
\patchcmd{\LT@caption}
  {\bgroup}
  {\bgroup\global\LTpatch@captiontrue}
  {}{}
\patchcmd{\longtable}
  {\par}
  {\par\global\LTpatch@captionfalse}
  {}{}
\apptocmd{\endlongtable}
  {\ifLTpatch@caption\else\addtocounter{table}{-1}\fi}
  {}{}
\newif\ifLTpatch@caption
\makeatother

\begin{document}

\maketitle


\hypertarget{toc}{}
\tableofcontents
\newpage
\section[Introduction]{Guía completa de Git y GitHub desde cero}

\setcounter{secnumdepth}{3}

\setlength\LTleft{0pt}




Git es un sistema de control de versiones distribuido rápido, escalable
y eficiente. Fue creado por Linus Torvalds en 2005 para el desarrollo
del kernel de Linux.

\textbf{Características principales:} - Sistema distribuido (cada
desarrollador tiene una copia completa del historial) - Extremadamente
rápido - Soporte robusto para desarrollo no lineal (branching y merging)
- Integridad de datos garantizada mediante SHA-1

\subsection{Arquitectura de Git}\label{arquitectura-de-git}

Git almacena los datos como instantáneas (snapshots) del proyecto
completo, no como diferencias entre archivos. Cada commit es una
instantánea completa del estado del proyecto.

\textbf{Tres áreas principales:} 1. \textbf{Working Directory}: Tu
directorio de trabajo actual 2. \textbf{Staging Area (Index)}: Área de
preparación para el próximo commit 3. \textbf{Repository (.git
directory)}: Base de datos de objetos y metadata

\section{Instalación y
Configuración}\label{instalaciuxf3n-y-configuraciuxf3n}

\subsection{Instalación}\label{instalaciuxf3n}

\subsubsection{Método 1: Paquetes predeterminados (rápido y
estable)}\label{muxe9todo-1-paquetes-predeterminados-ruxe1pido-y-estable}

\textbf{Linux (Ubuntu/Debian):}

\begin{Shaded}
\begin{Highlighting}[]
\FunctionTok{sudo}\NormalTok{ apt{-}get update}
\FunctionTok{sudo}\NormalTok{ apt{-}get install git}
\end{Highlighting}
\end{Shaded}

\textbf{Linux (Fedora):}

\begin{Shaded}
\begin{Highlighting}[]
\FunctionTok{sudo}\NormalTok{ dnf install git}
\end{Highlighting}
\end{Shaded}

\textbf{macOS:}

\begin{Shaded}
\begin{Highlighting}[]
\ExtensionTok{brew}\NormalTok{ install git}
\end{Highlighting}
\end{Shaded}

\textbf{Windows:} Descargar desde: https://git-scm.com/download/win

\subsection{Verificar Instalación}\label{verificar-instalaciuxf3n}

\begin{Shaded}
\begin{Highlighting}[]
\FunctionTok{git} \AttributeTok{{-}{-}version}
\end{Highlighting}
\end{Shaded}

\subsubsection{Método 2: Desde la fuente (versión más
reciente)}\label{muxe9todo-2-desde-la-fuente-versiuxf3n-muxe1s-reciente}

\begin{enumerate}
\def\labelenumi{\arabic{enumi}.}
\item
  Instala las dependencias:

\begin{Shaded}
\begin{Highlighting}[]
\FunctionTok{sudo}\NormalTok{ apt update}
\FunctionTok{sudo}\NormalTok{ apt install libz{-}dev libssl{-}dev libcurl4{-}gnutls{-}dev libexpat1{-}dev gettext cmake gcc}
\end{Highlighting}
\end{Shaded}
\item
  Descarga y descomprime la versión deseada (ejemplo: 2.34.1):

\begin{Shaded}
\begin{Highlighting}[]
\FunctionTok{mkdir}\NormalTok{ tmp }\KeywordTok{\&\&} \BuiltInTok{cd}\NormalTok{ tmp}
\ExtensionTok{curl} \AttributeTok{{-}o}\NormalTok{ git.tar.gz https://mirrors.edge.kernel.org/pub/software/scm/git/git{-}2.34.1.tar.gz}
\FunctionTok{tar} \AttributeTok{{-}zxf}\NormalTok{ git.tar.gz}
\BuiltInTok{cd}\NormalTok{ git{-}}\PreprocessorTok{*}
\end{Highlighting}
\end{Shaded}
\item
  Compila e instala:

\begin{Shaded}
\begin{Highlighting}[]
\FunctionTok{make}\NormalTok{ prefix=/usr/local all}
\FunctionTok{sudo}\NormalTok{ make prefix=/usr/local install}
\BuiltInTok{exec}\NormalTok{ bash}
\end{Highlighting}
\end{Shaded}
\item
  Confirma la instalación:

\begin{Shaded}
\begin{Highlighting}[]
\FunctionTok{git} \AttributeTok{{-}{-}version}
\end{Highlighting}
\end{Shaded}
\end{enumerate}

\subsection{Configuración Inicial}\label{configuraciuxf3n-inicial}

\textbf{Configurar identidad (obligatorio):}

\begin{Shaded}
\begin{Highlighting}[]
\FunctionTok{git}\NormalTok{ config }\AttributeTok{{-}{-}global}\NormalTok{ user.name }\StringTok{"Tu Nombre"}
\FunctionTok{git}\NormalTok{ config }\AttributeTok{{-}{-}global}\NormalTok{ user.email }\StringTok{"tu.email@ejemplo.com"}
\end{Highlighting}
\end{Shaded}

\textbf{Editor por defecto:}

\begin{Shaded}
\begin{Highlighting}[]
\FunctionTok{git}\NormalTok{ config }\AttributeTok{{-}{-}global}\NormalTok{ core.editor }\StringTok{"vim"}
\FunctionTok{git}\NormalTok{ config }\AttributeTok{{-}{-}global}\NormalTok{ core.editor }\StringTok{"nano"}
\FunctionTok{git}\NormalTok{ config }\AttributeTok{{-}{-}global}\NormalTok{ core.editor }\StringTok{"code {-}{-}wait"}  \CommentTok{\# VS Code}
\end{Highlighting}
\end{Shaded}

\textbf{Configurar rama principal:}

\begin{Shaded}
\begin{Highlighting}[]
\FunctionTok{git}\NormalTok{ config }\AttributeTok{{-}{-}global}\NormalTok{ init.defaultBranch main}
\end{Highlighting}
\end{Shaded}

\textbf{Ver configuración:}

\begin{Shaded}
\begin{Highlighting}[]
\FunctionTok{git}\NormalTok{ config }\AttributeTok{{-}{-}list}
\FunctionTok{git}\NormalTok{ config }\AttributeTok{{-}{-}list} \AttributeTok{{-}{-}show{-}origin}  \CommentTok{\# Ver de dónde viene cada configuración}
\FunctionTok{git}\NormalTok{ config user.name             }\CommentTok{\# Ver configuración específica}
\end{Highlighting}
\end{Shaded}

\textbf{Niveles de configuración:}

\begin{itemize}
\tightlist
\item
  \texttt{-\/-system}: Para todos los usuarios del sistema
  (/etc/gitconfig)
\item
  \texttt{-\/-global}: Para tu usuario (\textasciitilde/.gitconfig)
\item
  \texttt{-\/-local}: Para el repositorio actual (.git/config)
  {[}predeterminado{]}
\end{itemize}

\subsection{Configuración de claves SSH para
GitHub}\label{configuraciuxf3n-de-claves-ssh-para-github}

\subsubsection{Generar una clave SSH}\label{generar-una-clave-ssh}

\begin{enumerate}
\def\labelenumi{\arabic{enumi}.}
\item
  Verifica claves existentes:

\begin{Shaded}
\begin{Highlighting}[]
\FunctionTok{ls} \AttributeTok{{-}al}\NormalTok{ \textasciitilde{}/.ssh}
\end{Highlighting}
\end{Shaded}

  Si no hay claves, crea el directorio:
  \texttt{mkdir\ \textasciitilde{}/.ssh}.
\item
  Genera un par de claves:

\begin{Shaded}
\begin{Highlighting}[]
\FunctionTok{ssh{-}keygen} \AttributeTok{{-}t}\NormalTok{ rsa }\AttributeTok{{-}b}\NormalTok{ 4096 }\AttributeTok{{-}C} \StringTok{"tu.email@example.com"}
\end{Highlighting}
\end{Shaded}

  Acepta el nombre predeterminado y añade una contraseña (opcional).
\end{enumerate}

\subsubsection{Añadir la clave a
ssh-agent}\label{auxf1adir-la-clave-a-ssh-agent}

\begin{enumerate}
\def\labelenumi{\arabic{enumi}.}
\item
  Inicia el agente:

\begin{Shaded}
\begin{Highlighting}[]
\BuiltInTok{eval} \StringTok{"}\VariableTok{$(}\FunctionTok{ssh{-}agent} \AttributeTok{{-}s}\VariableTok{)}\StringTok{"}
\end{Highlighting}
\end{Shaded}
\item
  Añade la clave privada:

\begin{Shaded}
\begin{Highlighting}[]
\FunctionTok{ssh{-}add}\NormalTok{ \textasciitilde{}/.ssh/id\_rsa}
\end{Highlighting}
\end{Shaded}
\end{enumerate}

\subsubsection{Vincular la clave a
GitHub}\label{vincular-la-clave-a-github}

\begin{enumerate}
\def\labelenumi{\arabic{enumi}.}
\item
  Copia la clave pública:

  \begin{itemize}
  \tightlist
  \item
    Linux/Mac: \texttt{cat\ \textasciitilde{}/.ssh/id\_rsa.pub}
  \item
    Windows:
    \texttt{clip\ \textless{}\ \textasciitilde{}/.ssh/id\_rsa.pub}
  \end{itemize}
\item
  En GitHub, ve a \emph{Settings \textgreater{} SSH and GPG keys
  \textgreater{} New SSH key}, pega la clave y guárdala.
\item
  Prueba la conexión:

\begin{Shaded}
\begin{Highlighting}[]
\FunctionTok{ssh} \AttributeTok{{-}T}\NormalTok{ git@github.com}
\end{Highlighting}
\end{Shaded}

  Resultado esperado:
  \texttt{Hi\ tu\_usuario!\ You\textquotesingle{}ve\ successfully\ authenticated...}
\end{enumerate}

\section{Conceptos Fundamentales}\label{conceptos-fundamentales}

\subsection{Objetos de Git}\label{objetos-de-git}

Git maneja cuatro tipos de objetos:

\begin{enumerate}
\def\labelenumi{\arabic{enumi}.}
\tightlist
\item
  \textbf{Blob}: Contenido de archivos
\item
  \textbf{Tree}: Estructura de directorios (apunta a blobs y otros
  trees)
\item
  \textbf{Commit}: Instantánea del proyecto (apunta a un tree y commits
  padres)
\item
  \textbf{Tag}: Referencia nombrada a un commit específico
\end{enumerate}

\subsection{Referencias Simbólicas}\label{referencias-simbuxf3licas}

\begin{itemize}
\tightlist
\item
  \textbf{HEAD}: Apunta al commit actual (normalmente la punta de una
  rama)
\item
  \textbf{Rama (branch)}: Puntero móvil a un commit
\item
  \textbf{Tag}: Puntero fijo a un commit
\item
  \textbf{Remote}: Referencias a ramas en repositorios remotos
\end{itemize}

\subsection{Estados de Archivos}\label{estados-de-archivos}

\begin{enumerate}
\def\labelenumi{\arabic{enumi}.}
\tightlist
\item
  \textbf{Untracked}: Archivos nuevos que Git no rastrea
\item
  \textbf{Unmodified}: Archivos rastreados sin cambios
\item
  \textbf{Modified}: Archivos rastreados con cambios
\item
  \textbf{Staged}: Archivos preparados para el próximo commit
\end{enumerate}

\subsection{SHA-1 y Hashes}\label{sha-1-y-hashes}

Cada objeto en Git tiene un identificador único de 40 caracteres
hexadecimales (hash SHA-1). Puedes usar las primeras 7 caracteres para
referencias cortas.

\section{Comandos Básicos}\label{comandos-buxe1sicos}

\subsection{Crear e Inicializar
Repositorios}\label{crear-e-inicializar-repositorios}

\textbf{Crear nuevo repositorio:}

\begin{Shaded}
\begin{Highlighting}[]
\FunctionTok{git}\NormalTok{ init}
\FunctionTok{git}\NormalTok{ init nombre{-}proyecto}
\FunctionTok{git}\NormalTok{ init }\AttributeTok{{-}{-}bare}  \CommentTok{\# Repositorio sin working directory (para servidores)}
\end{Highlighting}
\end{Shaded}

\textbf{Clonar repositorio existente:}

\begin{Shaded}
\begin{Highlighting}[]
\FunctionTok{git}\NormalTok{ clone }\OperatorTok{\textless{}}\NormalTok{url}\OperatorTok{\textgreater{}}
\FunctionTok{git}\NormalTok{ clone https://github.com/usuario/repositorio.git}
\FunctionTok{git}\NormalTok{ clone }\OperatorTok{\textless{}}\NormalTok{url}\OperatorTok{\textgreater{}} \OperatorTok{\textless{}}\NormalTok{directorio{-}destino}\OperatorTok{\textgreater{}}
\end{Highlighting}
\end{Shaded}

Solo las últimas n confirmaciones:

\begin{Shaded}
\begin{Highlighting}[]
\FunctionTok{git}\NormalTok{ clone }\AttributeTok{{-}{-}depth}\NormalTok{ 1 }\OperatorTok{\textless{}}\NormalTok{url}\OperatorTok{\textgreater{}}  \CommentTok{\# Clone superficial (solo último commit)}
\FunctionTok{git}\NormalTok{ clone }\AttributeTok{{-}b} \OperatorTok{\textless{}}\NormalTok{rama}\OperatorTok{\textgreater{}} \OperatorTok{\textless{}}\NormalTok{url}\OperatorTok{\textgreater{}}  \CommentTok{\# Clonar rama específica, o,}
\FunctionTok{git}\NormalTok{ clone }\AttributeTok{{-}{-}branch}\OperatorTok{=}\NormalTok{nombre{-}rama }\OperatorTok{\textless{}}\NormalTok{url}\OperatorTok{\textgreater{}} 
\end{Highlighting}
\end{Shaded}

\subsection{Estado y Diferencias}\label{estado-y-diferencias}

\textbf{Ver estado del repositorio:}

\begin{Shaded}
\begin{Highlighting}[]
\FunctionTok{git}\NormalTok{ status}
\FunctionTok{git}\NormalTok{ status }\AttributeTok{{-}s}              \CommentTok{\# Formato corto}
\FunctionTok{git}\NormalTok{ status }\AttributeTok{{-}{-}short}         \CommentTok{\# Formato abreviado}
\FunctionTok{git}\NormalTok{ status }\AttributeTok{{-}b}              \CommentTok{\# Mostrar rama e información de tracking}
\end{Highlighting}
\end{Shaded}

\textbf{Ver cambios:}

\begin{Shaded}
\begin{Highlighting}[]
\FunctionTok{git}\NormalTok{ diff                   }\CommentTok{\# Cambios no staged (no confirmados)}
\FunctionTok{git}\NormalTok{ diff }\AttributeTok{{-}{-}staged}          \CommentTok{\# Cambios staged}
\FunctionTok{git}\NormalTok{ diff }\AttributeTok{{-}{-}cached}          \CommentTok{\# Sinónimo de {-}{-}staged}
\FunctionTok{git}\NormalTok{ diff HEAD              }\CommentTok{\# Todos los cambios desde último commit}
\FunctionTok{git}\NormalTok{ diff }\OperatorTok{\textless{}}\NormalTok{rama1}\OperatorTok{\textgreater{}} \OperatorTok{\textless{}}\NormalTok{rama2}\OperatorTok{\textgreater{}}   \CommentTok{\# Diferencias entre ramas}
\FunctionTok{git}\NormalTok{ diff }\OperatorTok{\textless{}}\NormalTok{commit1}\OperatorTok{\textgreater{}} \OperatorTok{\textless{}}\NormalTok{commit2}\OperatorTok{\textgreater{}}  \CommentTok{\# Diferencias entre commits}
\FunctionTok{git}\NormalTok{ diff }\AttributeTok{{-}{-}stat}            \CommentTok{\# Resumen estadístico}
\FunctionTok{git}\NormalTok{ diff }\AttributeTok{{-}{-}name{-}only}       \CommentTok{\# Solo nombres de archivos}
\end{Highlighting}
\end{Shaded}

\subsection{Añadir Archivos (Staging)}\label{auxf1adir-archivos-staging}

\begin{Shaded}
\begin{Highlighting}[]
\FunctionTok{git}\NormalTok{ add }\OperatorTok{\textless{}}\NormalTok{archivo}\OperatorTok{\textgreater{}}          \CommentTok{\# Añadir archivo específico}
\FunctionTok{git}\NormalTok{ add .                  }\CommentTok{\# Añadir todos los archivos del directorio actual}
\FunctionTok{git}\NormalTok{ add }\AttributeTok{{-}A}                 \CommentTok{\# Añadir todos los archivos del repositorio}
\FunctionTok{git}\NormalTok{ add }\PreprocessorTok{*}\NormalTok{.js               }\CommentTok{\# Añadir por patrón}
\FunctionTok{git}\NormalTok{ add }\AttributeTok{{-}p}                 \CommentTok{\# Añadir interactivamente por fragmentos}
\FunctionTok{git}\NormalTok{ add }\AttributeTok{{-}u}                 \CommentTok{\# Añadir solo archivos ya rastreados modificados}
\end{Highlighting}
\end{Shaded}

\subsection{Commits}\label{commits}

\textbf{Crear commit:}

\begin{Shaded}
\begin{Highlighting}[]
\FunctionTok{git}\NormalTok{ commit }\AttributeTok{{-}m} \StringTok{"Mensaje del commit"}
\FunctionTok{git}\NormalTok{ commit }\AttributeTok{{-}am} \StringTok{"Mensaje"}   \CommentTok{\# Add + commit (solo tracked files)}
\FunctionTok{git}\NormalTok{ commit                 }\CommentTok{\# Abre editor para mensaje largo}
\FunctionTok{git}\NormalTok{ commit }\AttributeTok{{-}{-}amend}         \CommentTok{\# Modificar último commit}
\FunctionTok{git}\NormalTok{ commit }\AttributeTok{{-}{-}amend} \AttributeTok{{-}m} \StringTok{"Nuevo mensaje"}  \CommentTok{\# Cambiar mensaje del último commit}
\FunctionTok{git}\NormalTok{ commit }\AttributeTok{{-}{-}amend} \AttributeTok{{-}{-}no{-}edit}  \CommentTok{\# Añadir cambios sin cambiar mensaje}
\end{Highlighting}
\end{Shaded}

\textbf{Buenas prácticas para mensajes:} - Primera línea: resumen
conciso (50 caracteres o menos) - Línea en blanco - Descripción
detallada (72 caracteres por línea) - Usar imperativo: ``Añade'' no
``Añadido''

\subsection{Historial}\label{historial}

\textbf{Ver historial:}

\begin{Shaded}
\begin{Highlighting}[]
\FunctionTok{git}\NormalTok{ log}
\FunctionTok{git}\NormalTok{ log }\AttributeTok{{-}{-}oneline}          \CommentTok{\# Una línea por commit}
\FunctionTok{git}\NormalTok{ log }\AttributeTok{{-}{-}graph}            \CommentTok{\# Vista gráfica}
\FunctionTok{git}\NormalTok{ log }\AttributeTok{{-}{-}all}              \CommentTok{\# Todas las ramas}
\FunctionTok{git}\NormalTok{ log }\AttributeTok{{-}{-}decorate}         \CommentTok{\# Mostrar referencias (ramas, tags)}
\FunctionTok{git}\NormalTok{ log }\AttributeTok{{-}p}                 \CommentTok{\# Mostrar diferencias (patch)}
\FunctionTok{git}\NormalTok{ log }\AttributeTok{{-}n}\NormalTok{ 5               }\CommentTok{\# Últimos 5 commits}
\FunctionTok{git}\NormalTok{ log }\AttributeTok{{-}{-}since}\OperatorTok{=}\StringTok{"2 weeks"}  \CommentTok{\# Commits de las últimas 2 semanas}
\FunctionTok{git}\NormalTok{ log }\AttributeTok{{-}{-}author}\OperatorTok{=}\StringTok{"Edison"}  \CommentTok{\# Commits de un autor}
\FunctionTok{git}\NormalTok{ log }\AttributeTok{{-}{-}grep}\OperatorTok{=}\StringTok{"palabra"}   \CommentTok{\# Buscar en mensajes de commit}
\FunctionTok{git}\NormalTok{ log }\AttributeTok{{-}{-}}\NormalTok{ archivo.txt     }\CommentTok{\# Historial de un archivo}
\FunctionTok{git}\NormalTok{ log }\AttributeTok{{-}{-}follow}\NormalTok{ archivo.txt  }\CommentTok{\# Incluir renombrados}
\end{Highlighting}
\end{Shaded}

\textbf{Formatos personalizados:}

\begin{Shaded}
\begin{Highlighting}[]
\FunctionTok{git}\NormalTok{ log }\AttributeTok{{-}{-}pretty}\OperatorTok{=}\NormalTok{format:}\StringTok{"\%h {-} \%an, \%ar : \%s"}
\FunctionTok{git}\NormalTok{ log }\AttributeTok{{-}{-}pretty}\OperatorTok{=}\NormalTok{format:}\StringTok{"\%h \%s"} \AttributeTok{{-}{-}graph}
\FunctionTok{git}\NormalTok{ log }\AttributeTok{{-}{-}graph} \AttributeTok{{-}{-}oneline} \AttributeTok{{-}{-}all} \AttributeTok{{-}{-}decorate}
\FunctionTok{git}\NormalTok{ config }\AttributeTok{{-}{-}global}\NormalTok{ alias.tree }\StringTok{"log {-}{-}graph {-}{-}decorate {-}{-}all {-}{-}oneline"}
\FunctionTok{git}\NormalTok{ tree}
\end{Highlighting}
\end{Shaded}

\textbf{Placeholders útiles:} - \texttt{\%h}: Hash abreviado -
\texttt{\%H}: Hash completo - \texttt{\%an}: Nombre del autor -
\texttt{\%ae}: Email del autor - \texttt{\%ad}: Fecha del autor -
\texttt{\%ar}: Fecha relativa - \texttt{\%s}: Mensaje del commit

\subsection{Deshacer Cambios}\label{deshacer-cambios}

\textbf{Descartar cambios en working directory:}

\begin{Shaded}
\begin{Highlighting}[]
\FunctionTok{git}\NormalTok{ checkout }\AttributeTok{{-}{-}} \OperatorTok{\textless{}}\NormalTok{archivo}\OperatorTok{\textgreater{}}  \CommentTok{\# Método antiguo}
\FunctionTok{git}\NormalTok{ restore }\OperatorTok{\textless{}}\NormalTok{archivo}\OperatorTok{\textgreater{}}      \CommentTok{\# Método moderno (Git 2.23+)}
\FunctionTok{git}\NormalTok{ restore .              }\CommentTok{\# Restaurar todos los archivos}
\end{Highlighting}
\end{Shaded}

\textbf{Quitar archivos del staging area:}

\begin{Shaded}
\begin{Highlighting}[]
\FunctionTok{git}\NormalTok{ reset HEAD }\OperatorTok{\textless{}}\NormalTok{archivo}\OperatorTok{\textgreater{}}   \CommentTok{\# Método antiguo}
\FunctionTok{git}\NormalTok{ restore }\AttributeTok{{-}{-}staged} \OperatorTok{\textless{}}\NormalTok{archivo}\OperatorTok{\textgreater{}}  \CommentTok{\# Método moderno}
\FunctionTok{git}\NormalTok{ reset HEAD             }\CommentTok{\# Quitar todos}
\end{Highlighting}
\end{Shaded}

\textbf{Deshacer commits:}

\begin{Shaded}
\begin{Highlighting}[]
\FunctionTok{git}\NormalTok{ reset }\AttributeTok{{-}{-}soft}\NormalTok{ HEAD\textasciitilde{}1    }\CommentTok{\# Mantiene cambios en staging}
\FunctionTok{git}\NormalTok{ reset }\AttributeTok{{-}{-}mixed}\NormalTok{ HEAD\textasciitilde{}1   }\CommentTok{\# Mantiene cambios en working directory [default]}
\FunctionTok{git}\NormalTok{ reset }\AttributeTok{{-}{-}hard}\NormalTok{ HEAD\textasciitilde{}1    }\CommentTok{\# CUIDADO: Elimina todo}
\FunctionTok{git}\NormalTok{ reset }\OperatorTok{\textless{}}\NormalTok{commit}\OperatorTok{\textgreater{}}         \CommentTok{\# Volver a un commit específico}
\end{Highlighting}
\end{Shaded}

\textbf{Revertir commit (seguro):}

\begin{Shaded}
\begin{Highlighting}[]
\FunctionTok{git}\NormalTok{ revert }\OperatorTok{\textless{}}\NormalTok{commit}\OperatorTok{\textgreater{}}        \CommentTok{\# Crea nuevo commit que deshace cambios}
\FunctionTok{git}\NormalTok{ revert HEAD            }\CommentTok{\# Revertir último commit}
\FunctionTok{git}\NormalTok{ revert }\AttributeTok{{-}{-}no{-}commit}\NormalTok{ HEAD\textasciitilde{}3..HEAD  }\CommentTok{\# Revertir últimos 3 commits}
\end{Highlighting}
\end{Shaded}

\subsection{Eliminar y Mover Archivos}\label{eliminar-y-mover-archivos}

\begin{Shaded}
\begin{Highlighting}[]
\FunctionTok{git}\NormalTok{ rm }\OperatorTok{\textless{}}\NormalTok{archivo}\OperatorTok{\textgreater{}}           \CommentTok{\# Eliminar archivo}
\FunctionTok{git}\NormalTok{ rm }\AttributeTok{{-}{-}cached} \OperatorTok{\textless{}}\NormalTok{archivo}\OperatorTok{\textgreater{}}  \CommentTok{\# Quitar del tracking pero mantener en disco}
\FunctionTok{git}\NormalTok{ rm }\AttributeTok{{-}r}\NormalTok{ directorio/      }\CommentTok{\# Eliminar directorio recursivamente}
\FunctionTok{git}\NormalTok{ mv }\OperatorTok{\textless{}}\NormalTok{origen}\OperatorTok{\textgreater{}} \OperatorTok{\textless{}}\NormalTok{destino}\OperatorTok{\textgreater{}}  \CommentTok{\# Renombrar/mover archivo}
\end{Highlighting}
\end{Shaded}

\section{Branching y Merging}\label{branching-y-merging}

\subsection{Conceptos de Ramas}\label{conceptos-de-ramas}

Una rama en Git es simplemente un puntero móvil a un commit. La rama por
defecto se llama \texttt{main} (anteriormente \texttt{master}).

\subsection{Operaciones con Ramas}\label{operaciones-con-ramas}

\textbf{Crear ramas:}

\begin{Shaded}
\begin{Highlighting}[]
\FunctionTok{git}\NormalTok{ branch }\OperatorTok{\textless{}}\NormalTok{nombre{-}rama}\OperatorTok{\textgreater{}}   \CommentTok{\# Crear rama}
\FunctionTok{git}\NormalTok{ checkout }\AttributeTok{{-}b} \OperatorTok{\textless{}}\NormalTok{nombre{-}rama}\OperatorTok{\textgreater{}}  \CommentTok{\# Crear y cambiar}
\FunctionTok{git}\NormalTok{ switch }\AttributeTok{{-}c} \OperatorTok{\textless{}}\NormalTok{nombre{-}rama}\OperatorTok{\textgreater{}}    \CommentTok{\# Método moderno (Git 2.23+)}
\end{Highlighting}
\end{Shaded}

\textbf{Listar ramas:}

\begin{Shaded}
\begin{Highlighting}[]
\FunctionTok{git}\NormalTok{ branch                 }\CommentTok{\# Ramas locales}
\FunctionTok{git}\NormalTok{ branch }\AttributeTok{{-}a}              \CommentTok{\# Todas las ramas (locales + remotas)}
\FunctionTok{git}\NormalTok{ branch }\AttributeTok{{-}r}              \CommentTok{\# Solo ramas remotas}
\FunctionTok{git}\NormalTok{ branch }\AttributeTok{{-}v}              \CommentTok{\# Con último commit de cada rama}
\FunctionTok{git}\NormalTok{ branch }\AttributeTok{{-}{-}merged}        \CommentTok{\# Ramas fusionadas con la actual}
\FunctionTok{git}\NormalTok{ branch }\AttributeTok{{-}{-}no{-}merged}     \CommentTok{\# Ramas no fusionadas}
\end{Highlighting}
\end{Shaded}

\textbf{Cambiar de rama:}

\begin{Shaded}
\begin{Highlighting}[]
\FunctionTok{git}\NormalTok{ checkout }\OperatorTok{\textless{}}\NormalTok{rama}\OperatorTok{\textgreater{}}        \CommentTok{\# Método antiguo}
\FunctionTok{git}\NormalTok{ switch }\OperatorTok{\textless{}}\NormalTok{rama}\OperatorTok{\textgreater{}}          \CommentTok{\# Método moderno (Git 2.23+)}
\FunctionTok{git}\NormalTok{ checkout }\AttributeTok{{-}}            \CommentTok{\# Volver a rama anterior}
\end{Highlighting}
\end{Shaded}

\textbf{Eliminar ramas:}

\begin{Shaded}
\begin{Highlighting}[]
\FunctionTok{git}\NormalTok{ branch }\AttributeTok{{-}d} \OperatorTok{\textless{}}\NormalTok{rama}\OperatorTok{\textgreater{}}       \CommentTok{\# Eliminar rama fusionada}
\FunctionTok{git}\NormalTok{ branch }\AttributeTok{{-}D} \OperatorTok{\textless{}}\NormalTok{rama}\OperatorTok{\textgreater{}}       \CommentTok{\# Forzar eliminación}
\FunctionTok{git}\NormalTok{ push origin }\AttributeTok{{-}{-}delete} \OperatorTok{\textless{}}\NormalTok{rama}\OperatorTok{\textgreater{}}  \CommentTok{\# Eliminar rama remota}
\end{Highlighting}
\end{Shaded}

\textbf{Renombrar rama:}

\begin{Shaded}
\begin{Highlighting}[]
\FunctionTok{git}\NormalTok{ branch }\AttributeTok{{-}m} \OperatorTok{\textless{}}\NormalTok{nuevo{-}nombre}\OperatorTok{\textgreater{}}     \CommentTok{\# Renombrar rama actual}
\FunctionTok{git}\NormalTok{ branch }\AttributeTok{{-}m} \OperatorTok{\textless{}}\NormalTok{viejo}\OperatorTok{\textgreater{}} \OperatorTok{\textless{}}\NormalTok{nuevo}\OperatorTok{\textgreater{}}    \CommentTok{\# Renombrar otra rama}
\end{Highlighting}
\end{Shaded}

\subsection{Merging (Fusionar)}\label{merging-fusionar}

\textbf{Merge básico:}

\begin{Shaded}
\begin{Highlighting}[]
\FunctionTok{git}\NormalTok{ merge }\OperatorTok{\textless{}}\NormalTok{rama}\OperatorTok{\textgreater{}}           \CommentTok{\# Fusionar rama en la actual}
\FunctionTok{git}\NormalTok{ merge }\AttributeTok{{-}{-}no{-}ff} \OperatorTok{\textless{}}\NormalTok{rama}\OperatorTok{\textgreater{}}   \CommentTok{\# Forzar commit de merge}
\FunctionTok{git}\NormalTok{ merge }\AttributeTok{{-}{-}squash} \OperatorTok{\textless{}}\NormalTok{rama}\OperatorTok{\textgreater{}}  \CommentTok{\# Fusionar como un solo commit}
\FunctionTok{git}\NormalTok{ merge }\AttributeTok{{-}{-}abort}          \CommentTok{\# Cancelar merge en conflicto}
\end{Highlighting}
\end{Shaded}

\textbf{Tipos de merge:} 1. \textbf{Fast-forward}: Avance directo sin
commit de merge 2. \textbf{Three-way merge}: Crea commit de merge con
dos padres 3. \textbf{Squash merge}: Combina todos los commits en uno

\subsection{Rebasing}\label{rebasing}

Rebase reescribe el historial moviendo commits a una nueva base.

\begin{Shaded}
\begin{Highlighting}[]
\FunctionTok{git}\NormalTok{ rebase }\OperatorTok{\textless{}}\NormalTok{rama{-}base}\OperatorTok{\textgreater{}}     \CommentTok{\# Rebasar rama actual}
\FunctionTok{git}\NormalTok{ rebase main            }\CommentTok{\# Rebasar sobre main}
\FunctionTok{git}\NormalTok{ rebase }\AttributeTok{{-}i}\NormalTok{ HEAD\textasciitilde{}3       }\CommentTok{\# Rebase interactivo (últimos 3 commits)}
\FunctionTok{git}\NormalTok{ rebase }\AttributeTok{{-}{-}continue}      \CommentTok{\# Continuar después de resolver conflictos}
\FunctionTok{git}\NormalTok{ rebase }\AttributeTok{{-}{-}abort}         \CommentTok{\# Cancelar rebase}
\FunctionTok{git}\NormalTok{ rebase }\AttributeTok{{-}{-}skip}          \CommentTok{\# Saltar commit actual}
\end{Highlighting}
\end{Shaded}

\textbf{Rebase interactivo - comandos:} - \texttt{pick}: Usar commit -
\texttt{reword}: Cambiar mensaje - \texttt{edit}: Editar commit -
\texttt{squash}: Fusionar con commit anterior - \texttt{fixup}: Como
squash pero descarta mensaje - \texttt{drop}: Eliminar commit

\textbf{⚠️ REGLA DE ORO: Nunca hagas rebase de commits
públicos/compartidos}

\subsection{Cherry-pick}\label{cherry-pick}

Aplicar commits específicos de una rama a otra:

\begin{Shaded}
\begin{Highlighting}[]
\FunctionTok{git}\NormalTok{ cherry{-}pick }\OperatorTok{\textless{}}\NormalTok{commit}\OperatorTok{\textgreater{}}   \CommentTok{\# Aplicar un commit}
\FunctionTok{git}\NormalTok{ cherry{-}pick }\OperatorTok{\textless{}}\NormalTok{commit1}\OperatorTok{\textgreater{}} \OperatorTok{\textless{}}\NormalTok{commit2}\OperatorTok{\textgreater{}}  \CommentTok{\# Varios commits}
\FunctionTok{git}\NormalTok{ cherry{-}pick }\OperatorTok{\textless{}}\NormalTok{commit1}\OperatorTok{\textgreater{}}\NormalTok{..}\OperatorTok{\textless{}}\NormalTok{commit2}\OperatorTok{\textgreater{}}  \CommentTok{\# Rango de commits}
\FunctionTok{git}\NormalTok{ cherry{-}pick }\AttributeTok{{-}{-}continue} \CommentTok{\# Continuar después de conflictos}
\FunctionTok{git}\NormalTok{ cherry{-}pick }\AttributeTok{{-}{-}abort}    \CommentTok{\# Cancelar}
\end{Highlighting}
\end{Shaded}

\section{Trabajo Remoto}\label{trabajo-remoto}

\subsection{Repositorios Remotos}\label{repositorios-remotos}

\textbf{Ver remotos:}

\begin{Shaded}
\begin{Highlighting}[]
\FunctionTok{git}\NormalTok{ remote                 }\CommentTok{\# Listar remotos}
\FunctionTok{git}\NormalTok{ remote }\AttributeTok{{-}v}              \CommentTok{\# Con URLs}
\FunctionTok{git}\NormalTok{ remote show origin     }\CommentTok{\# Información detallada}
\end{Highlighting}
\end{Shaded}

\textbf{Añadir remotos:}

\begin{Shaded}
\begin{Highlighting}[]
\FunctionTok{git}\NormalTok{ remote add }\OperatorTok{\textless{}}\NormalTok{nombre}\OperatorTok{\textgreater{}} \OperatorTok{\textless{}}\NormalTok{url}\OperatorTok{\textgreater{}}
\FunctionTok{git}\NormalTok{ remote add origin https://github.com/usuario/repo.git}
\end{Highlighting}
\end{Shaded}

\textbf{Modificar remotos:}

\begin{Shaded}
\begin{Highlighting}[]
\FunctionTok{git}\NormalTok{ remote rename }\OperatorTok{\textless{}}\NormalTok{viejo}\OperatorTok{\textgreater{}} \OperatorTok{\textless{}}\NormalTok{nuevo}\OperatorTok{\textgreater{}}
\FunctionTok{git}\NormalTok{ remote remove }\OperatorTok{\textless{}}\NormalTok{nombre}\OperatorTok{\textgreater{}}
\FunctionTok{git}\NormalTok{ remote set{-}url origin }\OperatorTok{\textless{}}\NormalTok{nueva{-}url}\OperatorTok{\textgreater{}}
\end{Highlighting}
\end{Shaded}

\subsection{Fetch, Pull y Push}\label{fetch-pull-y-push}

\textbf{Fetch (descargar sin fusionar):}

\begin{Shaded}
\begin{Highlighting}[]
\FunctionTok{git}\NormalTok{ fetch                  }\CommentTok{\# Fetch de origin}
\FunctionTok{git}\NormalTok{ fetch origin           }\CommentTok{\# Fetch de origin explícitamente}
\FunctionTok{git}\NormalTok{ fetch }\AttributeTok{{-}{-}all}            \CommentTok{\# Fetch de todos los remotos}
\FunctionTok{git}\NormalTok{ fetch origin }\OperatorTok{\textless{}}\NormalTok{rama}\OperatorTok{\textgreater{}}    \CommentTok{\# Fetch de rama específica}
\end{Highlighting}
\end{Shaded}

\textbf{Pull (fetch + merge):}

\begin{Shaded}
\begin{Highlighting}[]
\FunctionTok{git}\NormalTok{ pull                   }\CommentTok{\# Pull de rama actual}
\FunctionTok{git}\NormalTok{ pull origin main       }\CommentTok{\# Pull de rama específica}
\FunctionTok{git}\NormalTok{ pull }\AttributeTok{{-}{-}rebase}          \CommentTok{\# Pull con rebase en lugar de merge}
\FunctionTok{git}\NormalTok{ pull }\AttributeTok{{-}{-}rebase}\NormalTok{ origin main}
\end{Highlighting}
\end{Shaded}

\textbf{Push (enviar cambios):}

\begin{Shaded}
\begin{Highlighting}[]
\FunctionTok{git}\NormalTok{ push                   }\CommentTok{\# Push de rama actual}
\FunctionTok{git}\NormalTok{ push origin main       }\CommentTok{\# Push a rama específica}
\FunctionTok{git}\NormalTok{ push }\AttributeTok{{-}u}\NormalTok{ origin main    }\CommentTok{\# Push y establecer upstream}
\FunctionTok{git}\NormalTok{ push }\AttributeTok{{-}{-}all}             \CommentTok{\# Push de todas las ramas}
\FunctionTok{git}\NormalTok{ push }\AttributeTok{{-}{-}tags}            \CommentTok{\# Push de todos los tags}
\FunctionTok{git}\NormalTok{ push origin }\AttributeTok{{-}{-}delete} \OperatorTok{\textless{}}\NormalTok{rama}\OperatorTok{\textgreater{}}  \CommentTok{\# Eliminar rama remota}
\FunctionTok{git}\NormalTok{ push }\AttributeTok{{-}{-}force}           \CommentTok{\# ⚠️ Forzar push (PELIGROSO)}
\FunctionTok{git}\NormalTok{ push }\AttributeTok{{-}{-}force{-}with{-}lease}  \CommentTok{\# Forzar pero más seguro}
\end{Highlighting}
\end{Shaded}

\subsection{Tracking Branches}\label{tracking-branches}

\begin{Shaded}
\begin{Highlighting}[]
\FunctionTok{git}\NormalTok{ branch }\AttributeTok{{-}u}\NormalTok{ origin/main  }\CommentTok{\# Establecer upstream de rama actual}
\FunctionTok{git}\NormalTok{ branch }\AttributeTok{{-}{-}set{-}upstream{-}to}\OperatorTok{=}\NormalTok{origin/main}
\FunctionTok{git}\NormalTok{ branch }\AttributeTok{{-}vv}             \CommentTok{\# Ver tracking branches}
\end{Highlighting}
\end{Shaded}

\section{Comandos Avanzados}\label{comandos-avanzados}

\subsection{Stash (Guardar
Temporalmente)}\label{stash-guardar-temporalmente}

Guardar cambios sin hacer commit:

\begin{Shaded}
\begin{Highlighting}[]
\FunctionTok{git}\NormalTok{ stash                  }\CommentTok{\# Guardar cambios}
\FunctionTok{git}\NormalTok{ stash save }\StringTok{"mensaje"}   \CommentTok{\# Con mensaje descriptivo}
\FunctionTok{git}\NormalTok{ stash }\AttributeTok{{-}u}               \CommentTok{\# Incluir archivos untracked}
\FunctionTok{git}\NormalTok{ stash }\AttributeTok{{-}{-}all}            \CommentTok{\# Incluir todo (untracked + ignored)}
\FunctionTok{git}\NormalTok{ stash list             }\CommentTok{\# Listar stashes}
\FunctionTok{git}\NormalTok{ stash show             }\CommentTok{\# Ver cambios del último stash}
\FunctionTok{git}\NormalTok{ stash show }\AttributeTok{{-}p}          \CommentTok{\# Ver diff completo}
\FunctionTok{git}\NormalTok{ stash apply            }\CommentTok{\# Aplicar último stash (lo mantiene)}
\FunctionTok{git}\NormalTok{ stash pop              }\CommentTok{\# Aplicar y eliminar último stash}
\FunctionTok{git}\NormalTok{ stash apply stash@\{2\}  }\CommentTok{\# Aplicar stash específico}
\FunctionTok{git}\NormalTok{ stash drop stash@\{0\}   }\CommentTok{\# Eliminar stash específico}
\FunctionTok{git}\NormalTok{ stash clear            }\CommentTok{\# Eliminar todos los stashes}
\FunctionTok{git}\NormalTok{ stash branch }\OperatorTok{\textless{}}\NormalTok{rama}\OperatorTok{\textgreater{}}    \CommentTok{\# Crear rama desde stash}
\end{Highlighting}
\end{Shaded}

\subsection{Tags}\label{tags}

Marcar puntos importantes en el historial:

\textbf{Tags ligeros:}

\begin{Shaded}
\begin{Highlighting}[]
\FunctionTok{git}\NormalTok{ tag v1.0               }\CommentTok{\# Tag ligero}
\FunctionTok{git}\NormalTok{ tag v1.0 }\OperatorTok{\textless{}}\NormalTok{commit}\OperatorTok{\textgreater{}}      \CommentTok{\# Tag en commit específico}
\end{Highlighting}
\end{Shaded}

\textbf{Tags anotados (recomendado):}

\begin{Shaded}
\begin{Highlighting}[]
\FunctionTok{git}\NormalTok{ tag }\AttributeTok{{-}a}\NormalTok{ v1.0 }\AttributeTok{{-}m} \StringTok{"Versión 1.0"}
\FunctionTok{git}\NormalTok{ tag }\AttributeTok{{-}a}\NormalTok{ v1.0 }\OperatorTok{\textless{}}\NormalTok{commit}\OperatorTok{\textgreater{}}\NormalTok{ {-}m }\StringTok{"mensaje"}
\end{Highlighting}
\end{Shaded}

\textbf{Operaciones con tags:}

\begin{Shaded}
\begin{Highlighting}[]
\FunctionTok{git}\NormalTok{ tag                    }\CommentTok{\# Listar tags}
\FunctionTok{git}\NormalTok{ tag }\AttributeTok{{-}l} \StringTok{"v1.8.*"}        \CommentTok{\# Listar con patrón}
\FunctionTok{git}\NormalTok{ show v1.0              }\CommentTok{\# Ver información del tag}
\FunctionTok{git}\NormalTok{ push origin v1.0       }\CommentTok{\# Push de tag específico}
\FunctionTok{git}\NormalTok{ push origin }\AttributeTok{{-}{-}tags}     \CommentTok{\# Push de todos los tags}
\FunctionTok{git}\NormalTok{ tag }\AttributeTok{{-}d}\NormalTok{ v1.0            }\CommentTok{\# Eliminar tag local}
\FunctionTok{git}\NormalTok{ push origin }\AttributeTok{{-}{-}delete}\NormalTok{ v1.0  }\CommentTok{\# Eliminar tag remoto}
\FunctionTok{git}\NormalTok{ checkout v1.0          }\CommentTok{\# Checkout a tag (detached HEAD)}
\end{Highlighting}
\end{Shaded}

\subsection{Buscar y Encontrar}\label{buscar-y-encontrar}

\textbf{Grep (buscar en archivos):}

\begin{Shaded}
\begin{Highlighting}[]
\FunctionTok{git}\NormalTok{ grep }\StringTok{"texto"}           \CommentTok{\# Buscar en working directory}
\FunctionTok{git}\NormalTok{ grep }\StringTok{"texto"} \OperatorTok{\textless{}}\NormalTok{commit}\OperatorTok{\textgreater{}}  \CommentTok{\# Buscar en commit específico}
\FunctionTok{git}\NormalTok{ grep }\AttributeTok{{-}n} \StringTok{"texto"}        \CommentTok{\# Con números de línea}
\FunctionTok{git}\NormalTok{ grep }\AttributeTok{{-}{-}count} \StringTok{"texto"}   \CommentTok{\# Contar coincidencias}
\FunctionTok{git}\NormalTok{ grep }\AttributeTok{{-}i} \StringTok{"texto"}        \CommentTok{\# Case insensitive}
\end{Highlighting}
\end{Shaded}

\textbf{Bisect (búsqueda binaria de bugs):}

\begin{Shaded}
\begin{Highlighting}[]
\FunctionTok{git}\NormalTok{ bisect start           }\CommentTok{\# Iniciar bisect}
\FunctionTok{git}\NormalTok{ bisect bad             }\CommentTok{\# Marcar commit actual como malo}
\FunctionTok{git}\NormalTok{ bisect good }\OperatorTok{\textless{}}\NormalTok{commit}\OperatorTok{\textgreater{}}   \CommentTok{\# Marcar commit bueno conocido}
\CommentTok{\# Git checkout commits intermedios automáticamente}
\FunctionTok{git}\NormalTok{ bisect good            }\CommentTok{\# Marcar como bueno}
\FunctionTok{git}\NormalTok{ bisect bad             }\CommentTok{\# Marcar como malo}
\FunctionTok{git}\NormalTok{ bisect reset           }\CommentTok{\# Terminar bisect}
\end{Highlighting}
\end{Shaded}

\textbf{Blame (ver quién modificó cada línea):}

\begin{Shaded}
\begin{Highlighting}[]
\FunctionTok{git}\NormalTok{ blame }\OperatorTok{\textless{}}\NormalTok{archivo}\OperatorTok{\textgreater{}}        \CommentTok{\# Ver autor de cada línea}
\FunctionTok{git}\NormalTok{ blame }\AttributeTok{{-}L}\NormalTok{ 10,20 }\OperatorTok{\textless{}}\NormalTok{archivo}\OperatorTok{\textgreater{}}  \CommentTok{\# Solo líneas 10{-}20}
\FunctionTok{git}\NormalTok{ blame }\AttributeTok{{-}e} \OperatorTok{\textless{}}\NormalTok{archivo}\OperatorTok{\textgreater{}}     \CommentTok{\# Mostrar emails}
\FunctionTok{git}\NormalTok{ blame }\AttributeTok{{-}w} \OperatorTok{\textless{}}\NormalTok{archivo}\OperatorTok{\textgreater{}}     \CommentTok{\# Ignorar cambios de whitespace}
\end{Highlighting}
\end{Shaded}

\subsection{Reflog}\label{reflog}

Registro de todos los cambios a HEAD (incluso commits ``perdidos''):

\begin{Shaded}
\begin{Highlighting}[]
\FunctionTok{git}\NormalTok{ reflog                 }\CommentTok{\# Ver reflog}
\FunctionTok{git}\NormalTok{ reflog show            }\CommentTok{\# Igual que git reflog}
\FunctionTok{git}\NormalTok{ reflog }\AttributeTok{{-}{-}all}           \CommentTok{\# Reflog de todas las referencias}
\FunctionTok{git}\NormalTok{ reset }\AttributeTok{{-}{-}hard}\NormalTok{ HEAD@\{2\}  }\CommentTok{\# Volver a estado anterior}
\end{Highlighting}
\end{Shaded}

\subsection{Worktrees}\label{worktrees}

Tener múltiples working directories del mismo repositorio:

\begin{Shaded}
\begin{Highlighting}[]
\FunctionTok{git}\NormalTok{ worktree add }\OperatorTok{\textless{}}\NormalTok{path}\OperatorTok{\textgreater{}} \OperatorTok{\textless{}}\NormalTok{branch}\OperatorTok{\textgreater{}}  \CommentTok{\# Crear worktree}
\FunctionTok{git}\NormalTok{ worktree list          }\CommentTok{\# Listar worktrees}
\FunctionTok{git}\NormalTok{ worktree remove }\OperatorTok{\textless{}}\NormalTok{path}\OperatorTok{\textgreater{}} \CommentTok{\# Eliminar worktree}
\FunctionTok{git}\NormalTok{ worktree prune         }\CommentTok{\# Limpiar worktrees obsoletos}
\end{Highlighting}
\end{Shaded}

\subsection{Submodules}\label{submodules}

Incluir repositorios dentro de repositorios:

\begin{Shaded}
\begin{Highlighting}[]
\FunctionTok{git}\NormalTok{ submodule add }\OperatorTok{\textless{}}\NormalTok{url}\OperatorTok{\textgreater{}} \OperatorTok{\textless{}}\NormalTok{path}\OperatorTok{\textgreater{}}  \CommentTok{\# Añadir submodule}
\FunctionTok{git}\NormalTok{ submodule init         }\CommentTok{\# Inicializar submodules}
\FunctionTok{git}\NormalTok{ submodule update       }\CommentTok{\# Actualizar submodules}
\FunctionTok{git}\NormalTok{ submodule update }\AttributeTok{{-}{-}init} \AttributeTok{{-}{-}recursive}  \CommentTok{\# Init + update recursivo}
\FunctionTok{git}\NormalTok{ clone }\AttributeTok{{-}{-}recurse{-}submodules} \OperatorTok{\textless{}}\NormalTok{url}\OperatorTok{\textgreater{}}  \CommentTok{\# Clonar con submodules}
\end{Highlighting}
\end{Shaded}

\subsection{Archivos y Bundles}\label{archivos-y-bundles}

\textbf{Archive (crear archivo del repositorio):}

\begin{Shaded}
\begin{Highlighting}[]
\FunctionTok{git}\NormalTok{ archive }\AttributeTok{{-}{-}format}\OperatorTok{=}\NormalTok{zip }\AttributeTok{{-}{-}output}\OperatorTok{=}\NormalTok{proyecto.zip HEAD}
\FunctionTok{git}\NormalTok{ archive }\AttributeTok{{-}{-}format}\OperatorTok{=}\NormalTok{tar.gz }\AttributeTok{{-}{-}output}\OperatorTok{=}\NormalTok{proyecto.tar.gz main}
\end{Highlighting}
\end{Shaded}

\textbf{Bundle (repositorio portátil):}

\begin{Shaded}
\begin{Highlighting}[]
\FunctionTok{git}\NormalTok{ bundle create repo.bundle HEAD }\AttributeTok{{-}{-}all}
\FunctionTok{git}\NormalTok{ clone repo.bundle repo{-}clonado}
\FunctionTok{git}\NormalTok{ bundle verify repo.bundle}
\end{Highlighting}
\end{Shaded}

\section{Resolución de Conflictos}\label{resoluciuxf3n-de-conflictos}

\subsection{Identificar Conflictos}\label{identificar-conflictos}

Cuando hay conflictos durante merge o rebase:

\begin{Shaded}
\begin{Highlighting}[]
\FunctionTok{git}\NormalTok{ status                 }\CommentTok{\# Ver archivos con conflictos}
\FunctionTok{git}\NormalTok{ diff                   }\CommentTok{\# Ver conflictos}
\FunctionTok{git}\NormalTok{ diff }\AttributeTok{{-}{-}ours}            \CommentTok{\# Ver nuestra versión}
\FunctionTok{git}\NormalTok{ diff }\AttributeTok{{-}{-}theirs}          \CommentTok{\# Ver su versión}
\end{Highlighting}
\end{Shaded}

\subsection{Marcadores de Conflicto}\label{marcadores-de-conflicto}

\begin{verbatim}
<<<<<<< HEAD
Tu versión del código
=======
Su versión del código
>>>>>>> rama-a-fusionar
\end{verbatim}

\subsection{Resolver Conflictos}\label{resolver-conflictos}

\textbf{Manualmente:} 1. Editar archivos para resolver conflictos 2.
Eliminar marcadores
\texttt{\textless{}\textless{}\textless{}\textless{}\textless{}\textless{}\textless{}},
\texttt{=======},
\texttt{\textgreater{}\textgreater{}\textgreater{}\textgreater{}\textgreater{}\textgreater{}\textgreater{}}
3. \texttt{git\ add\ \textless{}archivo\textgreater{}} para marcar como
resuelto 4. \texttt{git\ commit} o \texttt{git\ rebase\ -\/-continue}

\textbf{Con herramientas:}

\begin{Shaded}
\begin{Highlighting}[]
\FunctionTok{git}\NormalTok{ mergetool              }\CommentTok{\# Abrir herramienta de merge}
\FunctionTok{git}\NormalTok{ mergetool }\AttributeTok{{-}{-}tool}\OperatorTok{=}\NormalTok{vimdiff}
\FunctionTok{git}\NormalTok{ config }\AttributeTok{{-}{-}global}\NormalTok{ merge.tool meld  }\CommentTok{\# Configurar herramienta}
\end{Highlighting}
\end{Shaded}

\textbf{Estrategias:}

\begin{Shaded}
\begin{Highlighting}[]
\FunctionTok{git}\NormalTok{ merge }\AttributeTok{{-}X}\NormalTok{ ours }\OperatorTok{\textless{}}\NormalTok{rama}\OperatorTok{\textgreater{}}   \CommentTok{\# Preferir nuestra versión}
\FunctionTok{git}\NormalTok{ merge }\AttributeTok{{-}X}\NormalTok{ theirs }\OperatorTok{\textless{}}\NormalTok{rama}\OperatorTok{\textgreater{}} \CommentTok{\# Preferir su versión}
\FunctionTok{git}\NormalTok{ checkout }\AttributeTok{{-}{-}ours} \OperatorTok{\textless{}}\NormalTok{archivo}\OperatorTok{\textgreater{}}    \CommentTok{\# Tomar nuestra versión}
\FunctionTok{git}\NormalTok{ checkout }\AttributeTok{{-}{-}theirs} \OperatorTok{\textless{}}\NormalTok{archivo}\OperatorTok{\textgreater{}}  \CommentTok{\# Tomar su versión}
\end{Highlighting}
\end{Shaded}

\section{Git Workflows}\label{git-workflows}

\subsection{Git Flow}\label{git-flow}

Modelo de branching con ramas específicas:

\begin{itemize}
\tightlist
\item
  \textbf{main}: Código de producción
\item
  \textbf{develop}: Rama de desarrollo
\item
  \textbf{feature/}*: Nuevas características
\item
  \textbf{release/}*: Preparación de releases
\item
  \textbf{hotfix/}*: Correcciones urgentes
\end{itemize}

\textbf{Comandos (con extensión git-flow):}

\begin{Shaded}
\begin{Highlighting}[]
\FunctionTok{git}\NormalTok{ flow init}
\FunctionTok{git}\NormalTok{ flow feature start nueva{-}caracteristica}
\FunctionTok{git}\NormalTok{ flow feature finish nueva{-}caracteristica}
\FunctionTok{git}\NormalTok{ flow release start 1.0.0}
\FunctionTok{git}\NormalTok{ flow release finish 1.0.0}
\FunctionTok{git}\NormalTok{ flow hotfix start fix{-}critico}
\FunctionTok{git}\NormalTok{ flow hotfix finish fix{-}critico}
\end{Highlighting}
\end{Shaded}

\subsection{GitHub Flow}\label{github-flow}

Workflow más simple: 1. Crear rama desde \texttt{main} 2. Hacer commits
3. Abrir Pull Request 4. Revisión de código 5. Merge a \texttt{main} 6.
Deploy automático

\subsection{Trunk-Based Development}\label{trunk-based-development}

\begin{itemize}
\tightlist
\item
  Una rama principal (\texttt{main} o \texttt{trunk})
\item
  Commits frecuentes y pequeños
\item
  Feature flags para ocultar trabajo en progreso
\item
  CI/CD robusto
\end{itemize}

\section{Ejemplos de uso diario de
Git}\label{ejemplos-de-uso-diario-de-git}

Flujos de trabajo más comunes que uso \textbf{todos los días} en
diferentes tipos de proyectos y entornos (individuales, equipos
pequeños, medianos y open-source).

\subsection{1. Flujo individual / Freelance / Proyectos personales (el
más frecuente para principiantes y side
projects)}\label{flujo-individual-freelance-proyectos-personales-el-muxe1s-frecuente-para-principiantes-y-side-projects}

\begin{Shaded}
\begin{Highlighting}[]
\CommentTok{\# Al empezar el día}
\FunctionTok{git}\NormalTok{ pull origin main                }\CommentTok{\# Siempre sincronizo primero}
\FunctionTok{git}\NormalTok{ status }\CommentTok{\# Muestra el estado actual del repositorio.}

\CommentTok{\# Trabajo en una nueva funcionalidad}
\FunctionTok{git}\NormalTok{ switch }\AttributeTok{{-}c}\NormalTok{ feat/login{-}social     }\CommentTok{\# Creo rama con convención clara}
\CommentTok{\# ... desarrollo varias horas ...}
\FunctionTok{git}\NormalTok{ add .}
\FunctionTok{git}\NormalTok{ commit }\AttributeTok{{-}m} \StringTok{"feat: implementar login con Google y GitHub"}

\CommentTok{\# Pequeña corrección rápida}
\FunctionTok{git}\NormalTok{ add src/components/LoginForm.tsx}
\FunctionTok{git}\NormalTok{ commit }\AttributeTok{{-}m} \StringTok{"fix: corregir validación de email en formulario"}

\CommentTok{\# Al final del día o antes de push}
\FunctionTok{git}\NormalTok{ branch }\AttributeTok{{-}M}\NormalTok{ main }\CommentTok{\# Muestra las ramas existentes.}
\FunctionTok{git}\NormalTok{ switch main}
\FunctionTok{git}\NormalTok{ pull                            }\CommentTok{\# Vuelvo a sincronizar por si alguien más empujó}
\FunctionTok{git}\NormalTok{ switch feat/login{-}social}
\FunctionTok{git}\NormalTok{ rebase main                     }\CommentTok{\# Mantengo mi rama actualizada (rebase limpio)}
\FunctionTok{git}\NormalTok{ switch main}
\FunctionTok{git}\NormalTok{ merge }\AttributeTok{{-}{-}ff{-}only}\NormalTok{ feat/login{-}social   }\CommentTok{\# Fast{-}forward si es posible}
\FunctionTok{git}\NormalTok{ push }\AttributeTok{{-}u}\NormalTok{ origin main }\CommentTok{\# Envía los cambios al repositorio remoto.}
\FunctionTok{git}\NormalTok{ branch }\AttributeTok{{-}d}\NormalTok{ feat/login{-}social     }\CommentTok{\# Elimino la rama ya que está integrada}

\CommentTok{\# Opcional: commit squash al final (muy común en proyectos pequeños)}
\FunctionTok{git}\NormalTok{ merge }\AttributeTok{{-}{-}squash}\NormalTok{ feat/login{-}social}
\FunctionTok{git}\NormalTok{ commit }\AttributeTok{{-}m} \StringTok{"feat: login con proveedores sociales + validaciones"}
\FunctionTok{git}\NormalTok{ push}
\end{Highlighting}
\end{Shaded}

\textbf{Ventajas}: Muy rápido, historial relativamente limpio, poco
overhead.

\subsection{2. Flujo típico de equipo mediano/pequeño con Pull Requests
(el más usado en
2025)}\label{flujo-tuxedpico-de-equipo-medianopequeuxf1o-con-pull-requests-el-muxe1s-usado-en-2025}

\begin{Shaded}
\begin{Highlighting}[]
\CommentTok{\# Inicio de jornada / después de stand{-}up}
\FunctionTok{git}\NormalTok{ fetch }\AttributeTok{{-}{-}prune}                   \CommentTok{\# Limpio referencias remotas obsoletas}
\FunctionTok{git}\NormalTok{ switch main}
\FunctionTok{git}\NormalTok{ pull}

\CommentTok{\# Nueva tarea (ej: ticket \#456 {-} mejorar rendimiento de dashboard)}
\FunctionTok{git}\NormalTok{ switch }\AttributeTok{{-}c}\NormalTok{ feature/456{-}dashboard{-}perf}

\CommentTok{\# Desarrollo normal (varios commits)}
\FunctionTok{git}\NormalTok{ commit }\AttributeTok{{-}m} \StringTok{"feat(dashboard): lazy loading de gráficos pesados"}
\FunctionTok{git}\NormalTok{ commit }\AttributeTok{{-}m} \StringTok{"refactor: extraer hook useChartData"}

\CommentTok{\# Antes de terminar el día}
\FunctionTok{git}\NormalTok{ rebase }\AttributeTok{{-}i}\NormalTok{ origin/main           }\CommentTok{\# O rebase main si ya está actualizado}
\CommentTok{\# Squash/reword si quiero limpiar commits antes de PR}

\CommentTok{\# Subo y creo Pull Request}
\FunctionTok{git}\NormalTok{ push }\AttributeTok{{-}u}\NormalTok{ origin feature/456{-}dashboard{-}perf}

\CommentTok{\# En GitHub/GitLab:}
\CommentTok{\#   → Creo PR → añado reviewers → espero revisión}
\CommentTok{\#   → Después de aprobar y pasar CI/CD:}
\CommentTok{\#     Merge squash / merge commit / rebase \& merge (depende de la política del equipo)}

\CommentTok{\# Una vez mergeado (normalmente por el reviewer o bot)}
\FunctionTok{git}\NormalTok{ switch main}
\FunctionTok{git}\NormalTok{ pull}
\FunctionTok{git}\NormalTok{ fetch }\AttributeTok{{-}{-}prune}
\FunctionTok{git}\NormalTok{ branch }\AttributeTok{{-}d}\NormalTok{ feature/456{-}dashboard{-}perf   }\CommentTok{\# Limpio localmente}
\end{Highlighting}
\end{Shaded}

\textbf{Variante muy común en 2025}:\\
Merge squash → un solo commit limpio en \texttt{main} con el título del
PR

\subsection{3. Corrección rápida de bug en producción (hotfix
diario)}\label{correcciuxf3n-ruxe1pida-de-bug-en-producciuxf3n-hotfix-diario}

\begin{Shaded}
\begin{Highlighting}[]
\CommentTok{\# Opción más rápida y segura (recomendada)}
\FunctionTok{git}\NormalTok{ switch main}
\FunctionTok{git}\NormalTok{ pull}
\FunctionTok{git}\NormalTok{ switch }\AttributeTok{{-}c}\NormalTok{ hotfix/789{-}boton{-}pago{-}falla}

\CommentTok{\# Arreglo rápido (1{-}3 commits)}
\FunctionTok{git}\NormalTok{ commit }\AttributeTok{{-}m} \StringTok{"fix: corregir validación de tarjeta en checkout"}

\FunctionTok{git}\NormalTok{ push }\AttributeTok{{-}u}\NormalTok{ origin hotfix/789{-}boton{-}pago{-}falla}

\CommentTok{\# → Crear PR rápido hacia main}
\CommentTok{\# → Revisión exprés (o auto{-}merge si es crítico)}
\CommentTok{\# → Merge (preferiblemente squash o rebase)}

\CommentTok{\# Después del deploy}
\FunctionTok{git}\NormalTok{ switch main}
\FunctionTok{git}\NormalTok{ pull}
\FunctionTok{git}\NormalTok{ tag hotfix{-}2025{-}12{-}30{-}boton{-}pago   }\CommentTok{\# O v1.2.3{-}hotfix si usas semver}
\FunctionTok{git}\NormalTok{ push }\AttributeTok{{-}{-}tags}
\end{Highlighting}
\end{Shaded}

\subsection{4. Flujo cuando trabajas en varias tareas al mismo tiempo
(muy
común)}\label{flujo-cuando-trabajas-en-varias-tareas-al-mismo-tiempo-muy-comuxfan}

\begin{Shaded}
\begin{Highlighting}[]
\CommentTok{\# Estoy en medio de feature/cuenta{-}premium}
\FunctionTok{git}\NormalTok{ status}

\CommentTok{\# → Aparece bug urgente}
\FunctionTok{git}\NormalTok{ stash push }\AttributeTok{{-}m} \StringTok{"WIP premium checkout"}

\CommentTok{\# Arreglo rápido}
\FunctionTok{git}\NormalTok{ switch }\AttributeTok{{-}c}\NormalTok{ hotfix/791{-}api{-}timeout}
\CommentTok{\# ... fix ...}
\FunctionTok{git}\NormalTok{ commit }\AttributeTok{{-}m} \StringTok{"fix: aumentar timeout en llamada a /reports"}
\FunctionTok{git}\NormalTok{ push }\AttributeTok{{-}u}\NormalTok{ origin hotfix/791{-}api{-}timeout}
\CommentTok{\# → PR rápido → merge}

\CommentTok{\# Vuelvo a lo mío}
\FunctionTok{git}\NormalTok{ switch feature/cuenta{-}premium}
\FunctionTok{git}\NormalTok{ stash pop}
\CommentTok{\# Continúo donde estaba...}
\end{Highlighting}
\end{Shaded}

\subsection{5. Mini-resumen de comandos que más se usan en la práctica
diaria
(2025)}\label{mini-resumen-de-comandos-que-muxe1s-se-usan-en-la-pruxe1ctica-diaria-2025}

\begin{longtable}[]{@{}
  >{\raggedright\arraybackslash}p{(\linewidth - 4\tabcolsep) * \real{0.3737}}
  >{\raggedright\arraybackslash}p{(\linewidth - 4\tabcolsep) * \real{0.5051}}
  >{\raggedright\arraybackslash}p{(\linewidth - 4\tabcolsep) * \real{0.1212}}@{}}
\toprule\noalign{}
\begin{minipage}[b]{\linewidth}\raggedright
Acción
\end{minipage} & \begin{minipage}[b]{\linewidth}\raggedright
Comando más común en 2025
\end{minipage} & \begin{minipage}[b]{\linewidth}\raggedright
Frecuencia
\end{minipage} \\
\midrule\noalign{}
\endhead
\bottomrule\noalign{}
\endlastfoot
Sincronizar al empezar & \texttt{git\ pull} /
\texttt{git\ fetch\ -\/-prune\ \&\&\ git\ pull} & ★★★★★ \\
Crear rama de tarea & \texttt{git\ switch\ -c\ feat/ticket-xxx-nombre} &
★★★★★ \\
Commit atómico & \texttt{git\ commit\ -m\ "tipo:\ descripción\ corta"} &
★★★★★ \\
Actualizar rama antes de PR & \texttt{git\ rebase\ main} o
\texttt{git\ merge\ main} & ★★★★ \\
Subir rama nueva & \texttt{git\ push\ -u\ origin\ nombre-rama} & ★★★★ \\
Limpiar después de merge & \texttt{git\ branch\ -d\ rama-antigua} &
★★★★ \\
Guardar trabajo temporal & \texttt{git\ stash\ push\ -m\ "..."} /
\texttt{git\ stash\ pop} & ★★★ \\
Ver estado rápido & \texttt{git\ status\ -sb} o alias \texttt{gs} &
★★★★★ \\
Historial bonito &
\texttt{git\ log\ -\/-oneline\ -\/-graph\ -\/-decorate\ -\/-all} &
★★★★ \\
Comparar con main & \texttt{git\ diff\ main...} & ★★★ \\
\end{longtable}

\textbf{Recomendación final para tu día a día}

Adopta el hábito de:

\begin{enumerate}
\def\labelenumi{\arabic{enumi}.}
\tightlist
\item
  Siempre \texttt{pull} al empezar
\item
  Trabajar en ramas cortas y con nombres claros
\item
  Commits pequeños y descriptivos (convención \textbf{Conventional
  Commits} es la más usada actualmente)
\item
  Actualizar tu rama frecuentemente (\texttt{rebase} o
  \texttt{merge\ main})
\item
  Eliminar ramas una vez integradas
\end{enumerate}

\section{Mejores Prácticas}\label{mejores-pruxe1cticas}

\subsection{Commits}\label{commits-1}

\begin{enumerate}
\def\labelenumi{\arabic{enumi}.}
\tightlist
\item
  \textbf{Commits atómicos}: Un commit = un cambio lógico
\item
  \textbf{Mensajes descriptivos}: Explica el ``qué'' y el ``por qué''
\item
  \textbf{Commits frecuentes}: Mejor muchos commits pequeños que pocos
  grandes
\item
  \textbf{No commitear archivos generados}: Usar .gitignore
\item
  \textbf{Revisar antes de commitear}: \texttt{git\ diff\ -\/-staged}
\end{enumerate}

\subsection{Branching}\label{branching}

\begin{enumerate}
\def\labelenumi{\arabic{enumi}.}
\tightlist
\item
  \textbf{Nombres descriptivos}: \texttt{feature/nueva-funcionalidad},
  \texttt{bugfix/issue-123}
\item
  \textbf{Ramas de vida corta}: Fusionar frecuentemente
\item
  \textbf{Mantener ramas actualizadas}: Hacer merge/rebase regular de
  main
\item
  \textbf{Eliminar ramas fusionadas}: Mantener repositorio limpio
\item
  \textbf{Proteger rama principal}: Requerir pull requests y revisiones
\end{enumerate}

\subsection{Colaboración}\label{colaboraciuxf3n}

\begin{enumerate}
\def\labelenumi{\arabic{enumi}.}
\tightlist
\item
  \textbf{Pull antes de push}: Evitar conflictos
\item
  \textbf{Revisar código}: Pull requests con revisiones
\item
  \textbf{No reescribir historial público}: Evitar rebase/amend de
  commits pusheados
\item
  \textbf{Comunicación}: Documentar decisiones importantes
\item
  \textbf{CI/CD}: Automatizar tests y despliegues
\end{enumerate}

\subsection{.gitignore}\label{gitignore}

Ejemplos comunes:

\begin{Shaded}
\begin{Highlighting}[]
\NormalTok{\# Node.js}
\NormalTok{node\_modules/}
\NormalTok{npm{-}debug.log}
\NormalTok{.env}

\NormalTok{\# Python}
\NormalTok{\_\_pycache\_\_/}
\NormalTok{*.py[cod]}
\NormalTok{venv/}
\NormalTok{.pytest\_cache/}

\NormalTok{\# IDE}
\NormalTok{.vscode/}
\NormalTok{.idea/}
\NormalTok{*.swp}

\NormalTok{\# OS}
\NormalTok{.DS\_Store}
\NormalTok{Thumbs.db}

\NormalTok{\# Build}
\NormalTok{dist/}
\NormalTok{build/}
\NormalTok{*.log}
\end{Highlighting}
\end{Shaded}

\textbf{Comandos útiles:}

\begin{Shaded}
\begin{Highlighting}[]
\FunctionTok{git}\NormalTok{ config }\AttributeTok{{-}{-}global}\NormalTok{ core.excludesfile \textasciitilde{}/.gitignore\_global}
\FunctionTok{git}\NormalTok{ check{-}ignore }\AttributeTok{{-}v} \OperatorTok{\textless{}}\NormalTok{archivo}\OperatorTok{\textgreater{}}  \CommentTok{\# Ver por qué archivo es ignorado}
\end{Highlighting}
\end{Shaded}

\section{Comandos de Bajo Nivel
(Plumbing)}\label{comandos-de-bajo-nivel-plumbing}

Comandos internos de Git para operaciones avanzadas:

\subsection{Inspección de Objetos}\label{inspecciuxf3n-de-objetos}

\begin{Shaded}
\begin{Highlighting}[]
\FunctionTok{git}\NormalTok{ cat{-}file }\AttributeTok{{-}p} \OperatorTok{\textless{}}\NormalTok{hash}\OperatorTok{\textgreater{}}     \CommentTok{\# Ver contenido de objeto}
\FunctionTok{git}\NormalTok{ cat{-}file }\AttributeTok{{-}t} \OperatorTok{\textless{}}\NormalTok{hash}\OperatorTok{\textgreater{}}     \CommentTok{\# Ver tipo de objeto}
\FunctionTok{git}\NormalTok{ cat{-}file }\AttributeTok{{-}s} \OperatorTok{\textless{}}\NormalTok{hash}\OperatorTok{\textgreater{}}     \CommentTok{\# Ver tamaño de objeto}
\FunctionTok{git}\NormalTok{ ls{-}tree }\OperatorTok{\textless{}}\NormalTok{tree{-}hash}\OperatorTok{\textgreater{}}    \CommentTok{\# Listar contenido de tree}
\FunctionTok{git}\NormalTok{ ls{-}files               }\CommentTok{\# Listar archivos en index}
\FunctionTok{git}\NormalTok{ ls{-}files }\AttributeTok{{-}s}            \CommentTok{\# Con información detallada}
\FunctionTok{git}\NormalTok{ ls{-}remote }\OperatorTok{\textless{}}\NormalTok{remote}\OperatorTok{\textgreater{}}     \CommentTok{\# Listar referencias remotas}
\end{Highlighting}
\end{Shaded}

\subsection{Manipulación del Index}\label{manipulaciuxf3n-del-index}

\begin{Shaded}
\begin{Highlighting}[]
\FunctionTok{git}\NormalTok{ update{-}index }\AttributeTok{{-}{-}add} \OperatorTok{\textless{}}\NormalTok{archivo}\OperatorTok{\textgreater{}}  \CommentTok{\# Añadir al index}
\FunctionTok{git}\NormalTok{ update{-}index }\AttributeTok{{-}{-}remove} \OperatorTok{\textless{}}\NormalTok{archivo}\OperatorTok{\textgreater{}}  \CommentTok{\# Quitar del index}
\FunctionTok{git}\NormalTok{ read{-}tree }\OperatorTok{\textless{}}\NormalTok{tree}\OperatorTok{\textgreater{}}       \CommentTok{\# Leer tree al index}
\FunctionTok{git}\NormalTok{ write{-}tree             }\CommentTok{\# Crear tree desde index}
\end{Highlighting}
\end{Shaded}

\begin{Shaded}
\begin{Highlighting}[]
\FunctionTok{git}\NormalTok{ show{-}ref               }\CommentTok{\# Mostrar todas las referencias}
\FunctionTok{git}\NormalTok{ update{-}ref refs/heads/main }\OperatorTok{\textless{}}\NormalTok{commit}\OperatorTok{\textgreater{}}  \CommentTok{\# Actualizar referencia}
\FunctionTok{git}\NormalTok{ symbolic{-}ref HEAD      }\CommentTok{\# Ver a qué apunta HEAD}
\FunctionTok{git}\NormalTok{ for{-}each{-}ref           }\CommentTok{\# Iterar sobre referencias}
\end{Highlighting}
\end{Shaded}

\subsection{Objetos y Hashes}\label{objetos-y-hashes}

\begin{Shaded}
\begin{Highlighting}[]
\FunctionTok{git}\NormalTok{ hash{-}object }\OperatorTok{\textless{}}\NormalTok{archivo}\OperatorTok{\textgreater{}}  \CommentTok{\# Calcular hash de archivo}
\FunctionTok{git}\NormalTok{ hash{-}object }\AttributeTok{{-}w} \OperatorTok{\textless{}}\NormalTok{archivo}\OperatorTok{\textgreater{}}  \CommentTok{\# Escribir objeto}
\FunctionTok{git}\NormalTok{ rev{-}parse HEAD         }\CommentTok{\# Convertir referencia a hash}
\FunctionTok{git}\NormalTok{ rev{-}list HEAD          }\CommentTok{\# Listar commits alcanzables}
\end{Highlighting}
\end{Shaded}

\section{Hooks y Automatización}\label{hooks-y-automatizaciuxf3n}

\subsection{Git Hooks}\label{git-hooks}

Scripts que se ejecutan automáticamente en eventos de Git.

\textbf{Ubicación:} \texttt{.git/hooks/}

\textbf{Hooks comunes:}

\textbf{Client-side:} - \texttt{pre-commit}: Antes de crear commit -
\texttt{prepare-commit-msg}: Antes de abrir editor de commit -
\texttt{commit-msg}: Validar mensaje de commit - \texttt{post-commit}:
Después de crear commit - \texttt{pre-push}: Antes de hacer push

\textbf{Server-side:} - \texttt{pre-receive}: Antes de aceptar push -
\texttt{update}: Por cada rama actualizada - \texttt{post-receive}:
Después de aceptar push

\textbf{Ejemplo pre-commit (tests):}

\begin{Shaded}
\begin{Highlighting}[]
\CommentTok{\#!/bin/sh}
\CommentTok{\# .git/hooks/pre{-}commit}

\ExtensionTok{npm}\NormalTok{ test}
\ControlFlowTok{if} \BuiltInTok{[} \VariableTok{$?} \OtherTok{{-}ne}\NormalTok{ 0 }\BuiltInTok{]}\KeywordTok{;} \ControlFlowTok{then}
    \BuiltInTok{echo} \StringTok{"Tests fallaron. Commit cancelado."}
    \BuiltInTok{exit}\NormalTok{ 1}
\ControlFlowTok{fi}
\end{Highlighting}
\end{Shaded}

\textbf{Ejemplo commit-msg (validar formato):}

\begin{Shaded}
\begin{Highlighting}[]
\CommentTok{\#!/bin/sh}
\CommentTok{\# .git/hooks/commit{-}msg}

\VariableTok{commit\_msg}\OperatorTok{=}\VariableTok{$(}\FunctionTok{cat} \StringTok{"}\VariableTok{$1}\StringTok{"}\VariableTok{)}
\VariableTok{pattern}\OperatorTok{=}\StringTok{"\^{}(feat|fix|docs|style|refactor|test|chore): .+"}

\ControlFlowTok{if} \OtherTok{! }\BuiltInTok{echo} \StringTok{"}\VariableTok{$commit\_msg}\StringTok{"} \KeywordTok{|} \FunctionTok{grep} \AttributeTok{{-}qE} \StringTok{"}\VariableTok{$pattern}\StringTok{"}\KeywordTok{;} \ControlFlowTok{then}
    \BuiltInTok{echo} \StringTok{"Error: El mensaje debe seguir el formato:"}
    \BuiltInTok{echo} \StringTok{"(feat|fix|docs|style|refactor|test|chore): descripción"}
    \BuiltInTok{exit}\NormalTok{ 1}
\ControlFlowTok{fi}
\end{Highlighting}
\end{Shaded}

\subsection{Aliases}\label{aliases}

Crear comandos personalizados:

\begin{Shaded}
\begin{Highlighting}[]
\CommentTok{\# Configurar aliases}
\FunctionTok{git}\NormalTok{ config }\AttributeTok{{-}{-}global}\NormalTok{ alias.co checkout}
\FunctionTok{git}\NormalTok{ config }\AttributeTok{{-}{-}global}\NormalTok{ alias.br branch}
\FunctionTok{git}\NormalTok{ config }\AttributeTok{{-}{-}global}\NormalTok{ alias.ci commit}
\FunctionTok{git}\NormalTok{ config }\AttributeTok{{-}{-}global}\NormalTok{ alias.st status}
\FunctionTok{git}\NormalTok{ config }\AttributeTok{{-}{-}global}\NormalTok{ alias.unstage }\StringTok{\textquotesingle{}reset HEAD {-}{-}\textquotesingle{}}
\FunctionTok{git}\NormalTok{ config }\AttributeTok{{-}{-}global}\NormalTok{ alias.last }\StringTok{\textquotesingle{}log {-}1 HEAD\textquotesingle{}}
\FunctionTok{git}\NormalTok{ config }\AttributeTok{{-}{-}global}\NormalTok{ alias.visual }\StringTok{\textquotesingle{}log {-}{-}oneline {-}{-}graph {-}{-}all {-}{-}decorate\textquotesingle{}}
\FunctionTok{git}\NormalTok{ config }\AttributeTok{{-}{-}global}\NormalTok{ alias.amend }\StringTok{\textquotesingle{}commit {-}{-}amend {-}{-}no{-}edit\textquotesingle{}}

\CommentTok{\# Usar aliases}
\FunctionTok{git}\NormalTok{ co main}
\FunctionTok{git}\NormalTok{ visual}
\end{Highlighting}
\end{Shaded}

\textbf{Aliases útiles:}

\begin{Shaded}
\begin{Highlighting}[]
\FunctionTok{git}\NormalTok{ config }\AttributeTok{{-}{-}global}\NormalTok{ alias.lg }\StringTok{"log {-}{-}graph {-}{-}pretty=format:\textquotesingle{}\%Cred\%h\%Creset {-}\%C(yellow)\%d\%Creset \%s \%Cgreen(\%cr) \%C(bold blue)\textless{}\%an\textgreater{}\%Creset\textquotesingle{} {-}{-}abbrev{-}commit"}
\FunctionTok{git}\NormalTok{ config }\AttributeTok{{-}{-}global}\NormalTok{ alias.undo }\StringTok{\textquotesingle{}reset HEAD\textasciitilde{}1 {-}{-}mixed\textquotesingle{}}
\FunctionTok{git}\NormalTok{ config }\AttributeTok{{-}{-}global}\NormalTok{ alias.contributors }\StringTok{\textquotesingle{}shortlog {-}sn\textquotesingle{}}
\end{Highlighting}
\end{Shaded}

\section{Troubleshooting}\label{troubleshooting}

\subsection{Problemas Comunes}\label{problemas-comunes}

\textbf{Descartar todos los cambios locales:}

\begin{Shaded}
\begin{Highlighting}[]
\FunctionTok{git}\NormalTok{ reset }\AttributeTok{{-}{-}hard}\NormalTok{ HEAD}
\FunctionTok{git}\NormalTok{ clean }\AttributeTok{{-}fd}  \CommentTok{\# Eliminar archivos untracked}
\end{Highlighting}
\end{Shaded}

\textbf{Recuperar commit eliminado:}

\begin{Shaded}
\begin{Highlighting}[]
\FunctionTok{git}\NormalTok{ reflog  }\CommentTok{\# Encontrar hash del commit}
\FunctionTok{git}\NormalTok{ checkout }\OperatorTok{\textless{}}\NormalTok{hash}\OperatorTok{\textgreater{}}
\FunctionTok{git}\NormalTok{ branch recovery }\OperatorTok{\textless{}}\NormalTok{hash}\OperatorTok{\textgreater{}}  \CommentTok{\# Crear rama desde commit}
\end{Highlighting}
\end{Shaded}

\textbf{Cambiar último commit:}

\begin{Shaded}
\begin{Highlighting}[]
\FunctionTok{git}\NormalTok{ commit }\AttributeTok{{-}{-}amend}  \CommentTok{\# Modificar mensaje o añadir archivos}
\end{Highlighting}
\end{Shaded}

\textbf{Mover commits a otra rama:}

\begin{Shaded}
\begin{Highlighting}[]
\FunctionTok{git}\NormalTok{ checkout rama{-}correcta}
\FunctionTok{git}\NormalTok{ cherry{-}pick }\OperatorTok{\textless{}}\NormalTok{commit}\OperatorTok{\textgreater{}}
\FunctionTok{git}\NormalTok{ checkout rama{-}incorrecta}
\FunctionTok{git}\NormalTok{ reset }\AttributeTok{{-}{-}hard}\NormalTok{ HEAD\textasciitilde{}1  }\CommentTok{\# Eliminar de rama incorrecta}
\end{Highlighting}
\end{Shaded}

\textbf{Dividir commit grande:}

\begin{Shaded}
\begin{Highlighting}[]
\FunctionTok{git}\NormalTok{ reset HEAD\textasciitilde{}1  }\CommentTok{\# Deshacer commit pero mantener cambios}
\FunctionTok{git}\NormalTok{ add }\OperatorTok{\textless{}}\NormalTok{archivo1}\OperatorTok{\textgreater{}}
\FunctionTok{git}\NormalTok{ commit }\AttributeTok{{-}m} \StringTok{"Primera parte"}
\FunctionTok{git}\NormalTok{ add }\OperatorTok{\textless{}}\NormalTok{archivo2}\OperatorTok{\textgreater{}}
\FunctionTok{git}\NormalTok{ commit }\AttributeTok{{-}m} \StringTok{"Segunda parte"}
\end{Highlighting}
\end{Shaded}

\textbf{Sincronizar fork con upstream:}

\begin{Shaded}
\begin{Highlighting}[]
\FunctionTok{git}\NormalTok{ remote add upstream }\OperatorTok{\textless{}}\NormalTok{url{-}original}\OperatorTok{\textgreater{}}
\FunctionTok{git}\NormalTok{ fetch upstream}
\FunctionTok{git}\NormalTok{ checkout main}
\FunctionTok{git}\NormalTok{ merge upstream/main}
\FunctionTok{git}\NormalTok{ push origin main}
\end{Highlighting}
\end{Shaded}

\textbf{Resolver ``detached HEAD'':}

\begin{Shaded}
\begin{Highlighting}[]
\FunctionTok{git}\NormalTok{ branch temp{-}branch  }\CommentTok{\# Crear rama desde HEAD actual}
\FunctionTok{git}\NormalTok{ checkout main}
\FunctionTok{git}\NormalTok{ merge temp{-}branch}
\end{Highlighting}
\end{Shaded}

\textbf{Limpiar referencias obsoletas:}

\begin{Shaded}
\begin{Highlighting}[]
\FunctionTok{git}\NormalTok{ remote prune origin  }\CommentTok{\# Eliminar referencias remotas obsoletas}
\FunctionTok{git}\NormalTok{ fetch }\AttributeTok{{-}{-}prune}  \CommentTok{\# Fetch y prune simultáneamente}
\end{Highlighting}
\end{Shaded}

\textbf{Comprimir repositorio:}

\begin{Shaded}
\begin{Highlighting}[]
\FunctionTok{git}\NormalTok{ gc  }\CommentTok{\# Garbage collection}
\FunctionTok{git}\NormalTok{ gc }\AttributeTok{{-}{-}aggressive}  \CommentTok{\# Más agresivo}
\FunctionTok{git}\NormalTok{ prune  }\CommentTok{\# Eliminar objetos inalcanzables}
\end{Highlighting}
\end{Shaded}

\textbf{Ver archivos grandes:}

\begin{Shaded}
\begin{Highlighting}[]
\FunctionTok{git}\NormalTok{ rev{-}list }\AttributeTok{{-}{-}objects} \AttributeTok{{-}{-}all} \DataTypeTok{\textbackslash{}}
  \KeywordTok{|} \FunctionTok{git}\NormalTok{ cat{-}file }\AttributeTok{{-}{-}batch{-}check}\OperatorTok{=}\StringTok{\textquotesingle{}\%(objecttype) \%(objectname) \%(objectsize) \%(rest)\textquotesingle{}} \DataTypeTok{\textbackslash{}}
  \KeywordTok{|} \FunctionTok{sed} \AttributeTok{{-}n} \StringTok{\textquotesingle{}s/\^{}blob //p\textquotesingle{}} \DataTypeTok{\textbackslash{}}
  \KeywordTok{|} \FunctionTok{sort} \AttributeTok{{-}{-}numeric{-}sort} \AttributeTok{{-}{-}key}\OperatorTok{=}\NormalTok{2 }\DataTypeTok{\textbackslash{}}
  \KeywordTok{|} \FunctionTok{tail} \AttributeTok{{-}n}\NormalTok{ 10}
\end{Highlighting}
\end{Shaded}

\textbf{Eliminar archivo del historial:}

\begin{Shaded}
\begin{Highlighting}[]
\FunctionTok{git}\NormalTok{ filter{-}branch }\AttributeTok{{-}{-}tree{-}filter} \StringTok{\textquotesingle{}rm {-}f archivo{-}sensible\textquotesingle{}}\NormalTok{ HEAD}
\CommentTok{\# O usar git{-}filter{-}repo (más rápido)}
\FunctionTok{git}\NormalTok{ filter{-}repo }\AttributeTok{{-}{-}path}\NormalTok{ archivo{-}sensible }\AttributeTok{{-}{-}invert{-}paths}
\end{Highlighting}
\end{Shaded}

\section{Referencia Rápida}\label{referencia-ruxe1pida}

\subsection{Sintaxis de Referencias}\label{sintaxis-de-referencias}

\begin{Shaded}
\begin{Highlighting}[]
\ExtensionTok{HEAD}           \CommentTok{\# Commit actual}
\ExtensionTok{HEAD\^{}}          \CommentTok{\# Padre de HEAD (HEAD\textasciitilde{}1)}
\ExtensionTok{HEAD\^{}\^{}}         \CommentTok{\# Abuelo de HEAD (HEAD\textasciitilde{}2)}
\ExtensionTok{HEAD\textasciitilde{}3}         \CommentTok{\# 3 commits antes de HEAD}
\ExtensionTok{main\^{}2}         \CommentTok{\# Segundo padre de main (en merge)}
\OperatorTok{\textless{}}\NormalTok{branch}\OperatorTok{\textgreater{}}\NormalTok{@\{yesterday\}  }\CommentTok{\# Posición ayer}
\ExtensionTok{HEAD@\{5\}}       \CommentTok{\# 5 movimientos atrás en reflog}
\end{Highlighting}
\end{Shaded}

\subsection{Especificar Rangos}\label{especificar-rangos}

\begin{Shaded}
\begin{Highlighting}[]
\OperatorTok{\textless{}}\NormalTok{commit1}\OperatorTok{\textgreater{}}\NormalTok{..}\OperatorTok{\textless{}}\NormalTok{commit2}\OperatorTok{\textgreater{}}   \CommentTok{\# Commits alcanzables desde commit2 pero no desde commit1}
\OperatorTok{\textless{}}\NormalTok{commit1}\OperatorTok{\textgreater{}}\NormalTok{...}\OperatorTok{\textless{}}\NormalTok{commit2}\OperatorTok{\textgreater{}}  \CommentTok{\# Commits alcanzables desde cualquiera pero no desde ambos}
\OperatorTok{\textless{}}\NormalTok{branch}\OperatorTok{\textgreater{}}\NormalTok{\^{}@             }\CommentTok{\# Todos los padres de branch}
\OperatorTok{\textless{}}\NormalTok{commit}\OperatorTok{\textgreater{}}\NormalTok{\^{}!             }\CommentTok{\# El commit pero no sus padres}
\end{Highlighting}
\end{Shaded}

\subsection{Variables de Entorno}\label{variables-de-entorno}

\begin{Shaded}
\begin{Highlighting}[]
\ExtensionTok{GIT\_AUTHOR\_NAME}        \CommentTok{\# Nombre del autor}
\ExtensionTok{GIT\_AUTHOR\_EMAIL}       \CommentTok{\# Email del autor}
\ExtensionTok{GIT\_COMMITTER\_NAME}     \CommentTok{\# Nombre del committer}
\ExtensionTok{GIT\_COMMITTER\_EMAIL}    \CommentTok{\# Email del committer}
\ExtensionTok{GIT\_EDITOR}             \CommentTok{\# Editor predeterminado}
\ExtensionTok{GIT\_PAGER}              \CommentTok{\# Pager para output}
\ExtensionTok{GIT\_TRACE}              \CommentTok{\# Activar tracing}
\end{Highlighting}
\end{Shaded}

\subsection{Configuración Útil}\label{configuraciuxf3n-uxfatil}

\begin{Shaded}
\begin{Highlighting}[]
\CommentTok{\# Colorear output}
\FunctionTok{git}\NormalTok{ config }\AttributeTok{{-}{-}global}\NormalTok{ color.ui auto}

\CommentTok{\# Guardar credenciales}
\FunctionTok{git}\NormalTok{ config }\AttributeTok{{-}{-}global}\NormalTok{ credential.helper cache}
\FunctionTok{git}\NormalTok{ config }\AttributeTok{{-}{-}global}\NormalTok{ credential.helper }\StringTok{\textquotesingle{}cache {-}{-}timeout=3600\textquotesingle{}}

\CommentTok{\# Autocorrección de comandos}
\FunctionTok{git}\NormalTok{ config }\AttributeTok{{-}{-}global}\NormalTok{ help.autocorrect 10}

\CommentTok{\# Rebase por defecto al hacer pull}
\FunctionTok{git}\NormalTok{ config }\AttributeTok{{-}{-}global}\NormalTok{ pull.rebase true}

\CommentTok{\# Prune automático}
\FunctionTok{git}\NormalTok{ config }\AttributeTok{{-}{-}global}\NormalTok{ fetch.prune true}

\CommentTok{\# Diff mejorado}
\FunctionTok{git}\NormalTok{ config }\AttributeTok{{-}{-}global}\NormalTok{ diff.algorithm histogram}

\CommentTok{\# Rerere (reuse recorded resolution)}
\FunctionTok{git}\NormalTok{ config }\AttributeTok{{-}{-}global}\NormalTok{ rerere.enabled true}
\end{Highlighting}
\end{Shaded}

\subsection{Comandos de Emergencia}\label{comandos-de-emergencia}

\begin{Shaded}
\begin{Highlighting}[]
\CommentTok{\# Deshacer }\AlertTok{TODO}\CommentTok{ y volver limpio}
\FunctionTok{git}\NormalTok{ reset }\AttributeTok{{-}{-}hard}\NormalTok{ HEAD}
\FunctionTok{git}\NormalTok{ clean }\AttributeTok{{-}fd}

\CommentTok{\# Recuperar trabajo después de reset {-}{-}hard}
\FunctionTok{git}\NormalTok{ reflog}
\FunctionTok{git}\NormalTok{ reset }\AttributeTok{{-}{-}hard} \OperatorTok{\textless{}}\NormalTok{hash}\OperatorTok{\textgreater{}}

\CommentTok{\# Salir de cualquier operación en progreso}
\FunctionTok{git}\NormalTok{ merge }\AttributeTok{{-}{-}abort}
\FunctionTok{git}\NormalTok{ rebase }\AttributeTok{{-}{-}abort}
\FunctionTok{git}\NormalTok{ cherry{-}pick }\AttributeTok{{-}{-}abort}

\CommentTok{\# Verificar integridad del repositorio}
\FunctionTok{git}\NormalTok{ fsck }\AttributeTok{{-}{-}full}

\CommentTok{\# Reparar repositorio corrupto}
\FunctionTok{git}\NormalTok{ fsck }\AttributeTok{{-}{-}full} \AttributeTok{{-}{-}no{-}dangling}
\FunctionTok{git}\NormalTok{ gc }\AttributeTok{{-}{-}aggressive} \AttributeTok{{-}{-}prune}\OperatorTok{=}\NormalTok{now}
\end{Highlighting}
\end{Shaded}

\subsection{Atajos de Teclado (Shell)}\label{atajos-de-teclado-shell}

Añadir a \texttt{\textasciitilde{}/.bashrc} o
\texttt{\textasciitilde{}/.zshrc}:

\begin{Shaded}
\begin{Highlighting}[]
\BuiltInTok{alias}\NormalTok{ g=}\StringTok{\textquotesingle{}git\textquotesingle{}}
\BuiltInTok{alias}\NormalTok{ gs=}\StringTok{\textquotesingle{}git status\textquotesingle{}}
\BuiltInTok{alias}\NormalTok{ ga=}\StringTok{\textquotesingle{}git add\textquotesingle{}}
\BuiltInTok{alias}\NormalTok{ gc=}\StringTok{\textquotesingle{}git commit\textquotesingle{}}
\BuiltInTok{alias}\NormalTok{ gp=}\StringTok{\textquotesingle{}git push\textquotesingle{}}
\BuiltInTok{alias}\NormalTok{ gl=}\StringTok{\textquotesingle{}git pull\textquotesingle{}}
\BuiltInTok{alias}\NormalTok{ gd=}\StringTok{\textquotesingle{}git diff\textquotesingle{}}
\BuiltInTok{alias}\NormalTok{ gco=}\StringTok{\textquotesingle{}git checkout\textquotesingle{}}
\BuiltInTok{alias}\NormalTok{ gb=}\StringTok{\textquotesingle{}git branch\textquotesingle{}}
\BuiltInTok{alias}\NormalTok{ glg=}\StringTok{\textquotesingle{}git log {-}{-}graph {-}{-}oneline {-}{-}all {-}{-}decorate\textquotesingle{}}
\end{Highlighting}
\end{Shaded}

\subsection{Formato de Commit
Convencional}\label{formato-de-commit-convencional}

\begin{verbatim}
<tipo>(<ámbito>): <descripción>

<cuerpo>

<pie>
\end{verbatim}

\textbf{Tipos:} - \texttt{feat}: Nueva funcionalidad - \texttt{fix}:
Corrección de bug - \texttt{docs}: Documentación - \texttt{style}:
Formato (no afecta código) - \texttt{refactor}: Refactorización -
\texttt{test}: Tests - \texttt{chore}: Tareas de mantenimiento -
\texttt{perf}: Mejora de rendimiento

\textbf{Ejemplo:}

\begin{verbatim}
feat(auth): añadir autenticación con OAuth2

Implementa login con Google y GitHub usando OAuth2.
Incluye manejo de tokens y refresh automático.

Closes #123
\end{verbatim}

\section{Recursos Adicionales}\label{recursos-adicionales}

\subsection{Documentación Oficial}\label{documentaciuxf3n-oficial}

\begin{itemize}
\tightlist
\item
  \textbf{Manual de Git}: \texttt{man\ git} o https://git-scm.com/docs
\item
  \textbf{Libro Pro Git}: https://git-scm.com/book/es/v2
\item
  \textbf{Git Reference}: https://git-scm.com/docs
\end{itemize}

\subsection{Tutoriales y Cursos}\label{tutoriales-y-cursos}

\begin{itemize}
\tightlist
\item
  \textbf{Learn Git Branching}: https://learngitbranching.js.org
\item
  \textbf{Git Tutorial de Atlassian}:
  https://www.atlassian.com/git/tutorials
\item
  \textbf{GitHub Learning Lab}: https://lab.github.com
\end{itemize}

\subsection{Herramientas}\label{herramientas}

\begin{itemize}
\tightlist
\item
  \textbf{GitKraken}: Cliente visual multiplataforma
\item
  \textbf{SourceTree}: Cliente visual de Atlassian
\item
  \textbf{Fork}: Cliente Git para Mac y Windows
\item
  \textbf{lazygit}: Cliente TUI (terminal) simple y potente
\item
  \textbf{tig}: Navegador de texto para repositorios Git
\end{itemize}

\subsection{Extensiones}\label{extensiones}

\begin{itemize}
\tightlist
\item
  \textbf{git-extras}: Comandos útiles adicionales
\item
  \textbf{git-flow}: Extensión para Git Flow workflow
\item
  \textbf{git-lfs}: Large File Storage
\item
  \textbf{git-filter-repo}: Reescritura avanzada de historial
\end{itemize}

\section{Glosario}\label{glosario}

\textbf{Blob}: Objeto que contiene el contenido de un archivo.

\textbf{Branch (Rama)}: Puntero móvil a un commit que representa una
línea de desarrollo.

\textbf{Checkout}: Cambiar el working directory a un commit, rama o tag
específico.

\textbf{Clone}: Crear una copia local de un repositorio remoto.

\textbf{Commit}: Instantánea del proyecto en un momento específico.

\textbf{Conflict (Conflicto)}: Situación donde Git no puede fusionar
cambios automáticamente.

\textbf{Detached HEAD}: Estado donde HEAD apunta directamente a un
commit en lugar de a una rama.

\textbf{Diff}: Diferencias entre dos versiones de archivos.

\textbf{Fast-forward}: Tipo de merge donde simplemente se avanza el
puntero de la rama.

\textbf{Fetch}: Descargar objetos y referencias desde un repositorio
remoto sin fusionar.

\textbf{Fork}: Copia de un repositorio en tu propia cuenta.

\textbf{HEAD}: Referencia al commit actual.

\textbf{Index}: Área de staging donde se preparan cambios para el
próximo commit.

\textbf{Merge}: Fusionar cambios de una rama a otra.

\textbf{Origin}: Nombre por defecto del repositorio remoto principal.

\textbf{Pull}: Fetch + Merge en un solo comando.

\textbf{Pull Request}: Solicitud para fusionar cambios (terminología de
GitHub).

\textbf{Push}: Enviar commits locales a un repositorio remoto.

\textbf{Rebase}: Mover o combinar commits a una nueva base.

\textbf{Remote}: Repositorio alojado en otro lugar (servidor).

\textbf{Repository (Repositorio)}: Colección de commits, ramas y
configuración de un proyecto.

\textbf{SHA-1}: Algoritmo hash usado para identificar objetos de Git.

\textbf{Staging Area}: Sinónimo de Index.

\textbf{Stash}: Guardar temporalmente cambios sin hacer commit.

\textbf{Tag}: Referencia permanente a un commit específico.

\textbf{Tree}: Objeto que representa un directorio.

\textbf{Upstream}: Rama remota que una rama local rastrea.

\textbf{Working Directory (Directorio de Trabajo)}: Directorio actual
con archivos del proyecto.

\section{Conclusión}\label{conclusiuxf3n}

Git es una herramienta poderosa y flexible que requiere práctica para
dominar. Esta guía cubre desde conceptos básicos hasta técnicas
avanzadas, pero la mejor forma de aprender es usándolo en proyectos
reales.

\textbf{Consejos finales:}

\begin{enumerate}
\def\labelenumi{\arabic{enumi}.}
\tightlist
\item
  Practica regularmente
\item
  Experimenta en repositorios de prueba
\item
  Lee los mensajes de error (son informativos)
\item
  Usa \texttt{git\ help\ \textless{}comando\textgreater{}} cuando tengas
  dudas
\item
  No temas equivocarte (casi todo se puede deshacer)
\end{enumerate}

\textbf{Comandos más importantes para recordar:} -
\texttt{echo\ "\#\ Léeme"\ \textgreater{}\textgreater{}\ README.md}:
Crea un archivo README.md con el texto ``\# Léeme''. -
\texttt{git\ init}: Inicia un nuevo repositorio en el directorio actual.
- \texttt{git\ status} - Siempre saber dónde estás - \texttt{git\ add} -
Preparar cambios - \texttt{git\ commit\ -m\ "Primer\ commit"} - Guardar
cambios - \texttt{git\ branch\ -M\ main}: Muestra las ramas existentes.
- \texttt{git\ push\ -u\ origin\ main} - Compartir cambios -
\texttt{git\ pull} - Obtener cambios - \texttt{git\ log} - Ver historial
de commits - \texttt{git\ diff} - Ver diferencias -
\texttt{git\ checkout\ {[}branch{]}}: Cambia a una rama específica. -
\texttt{git\ merge\ {[}branch{]}}: Fusiona una rama con la actual.

\section{Publicaciones Similares}\label{publicaciones-similares}

Si te interesó este artículo, te recomendamos que explores otros blogs y
recursos relacionados que pueden ampliar tus conocimientos. Aquí te dejo
algunas sugerencias:

\begin{enumerate}
\def\labelenumi{\arabic{enumi}.}
\tightlist
\item
  \href{https://chaska-x.netlify.app/operating-system/2017-05-21-comandos-de-informacion-windows/index.pdf}{\faIcon{file-pdf}}
  \href{https://chaska-x.netlify.app/operating-system/2017-05-21-comandos-de-informacion-windows}{Comandos
  De Informacion Windows}
\item
  \href{https://chaska-x.netlify.app/operating-system/2019-06-19-adb/index.pdf}{\faIcon{file-pdf}}
  \href{https://chaska-x.netlify.app/operating-system/2019-06-19-adb}{Adb}
\item
  \href{https://chaska-x.netlify.app/operating-system/2021-08-17-limpieza-y-optimizacion-de-pc/index.pdf}{\faIcon{file-pdf}}
  \href{https://chaska-x.netlify.app/operating-system/2021-08-17-limpieza-y-optimizacion-de-pc}{Limpieza
  Y Optimizacion De Pc}
\item
  \href{https://chaska-x.netlify.app/operating-system/2021-10-21-usando-apk-en-windown-11/index.pdf}{\faIcon{file-pdf}}
  \href{https://chaska-x.netlify.app/operating-system/2021-10-21-usando-apk-en-windown-11}{Usando
  Apk En Windown 11}
\item
  \href{https://chaska-x.netlify.app/operating-system/2022-05-12-gestionar-versiones-de-jdk-en-kubuntu/index.pdf}{\faIcon{file-pdf}}
  \href{https://chaska-x.netlify.app/operating-system/2022-05-12-gestionar-versiones-de-jdk-en-kubuntu}{Gestionar
  Versiones De Jdk En Kubuntu}
\item
  \href{https://chaska-x.netlify.app/operating-system/2022-07-21-instalar-tor-browser/index.pdf}{\faIcon{file-pdf}}
  \href{https://chaska-x.netlify.app/operating-system/2022-07-21-instalar-tor-browser}{Instalar
  Tor Browser}
\item
  \href{https://chaska-x.netlify.app/operating-system/2022-08-14-crear-enlaces-duros-o-hard-link-en-linux/index.pdf}{\faIcon{file-pdf}}
  \href{https://chaska-x.netlify.app/operating-system/2022-08-14-crear-enlaces-duros-o-hard-link-en-linux}{Crear
  Enlaces Duros O Hard Link En Linux}
\item
  \href{https://chaska-x.netlify.app/operating-system/2022-09-27-comandos-vim/index.pdf}{\faIcon{file-pdf}}
  \href{https://chaska-x.netlify.app/operating-system/2022-09-27-comandos-vim}{Comandos
  Vim}
\item
  \href{https://chaska-x.netlify.app/operating-system/2023-02-16-guia-de-git-y-github/index.pdf}{\faIcon{file-pdf}}
  \href{https://chaska-x.netlify.app/operating-system/2023-02-16-guia-de-git-y-github}{Guia
  De Git Y Github}
\item
  \href{https://chaska-x.netlify.app/operating-system/2023-05-02-00-primeros-pasos-en-linux/index.pdf}{\faIcon{file-pdf}}
  \href{https://chaska-x.netlify.app/operating-system/2023-05-02-00-primeros-pasos-en-linux}{00
  Primeros Pasos En Linux}
\item
  \href{https://chaska-x.netlify.app/operating-system/2023-06-17-01-introduccion-linux/index.pdf}{\faIcon{file-pdf}}
  \href{https://chaska-x.netlify.app/operating-system/2023-06-17-01-introduccion-linux}{01
  Introduccion Linux}
\item
  \href{https://chaska-x.netlify.app/operating-system/2023-06-18-02-distribuciones-linux/index.pdf}{\faIcon{file-pdf}}
  \href{https://chaska-x.netlify.app/operating-system/2023-06-18-02-distribuciones-linux}{02
  Distribuciones Linux}
\item
  \href{https://chaska-x.netlify.app/operating-system/2023-06-19-03-instalacion-linux/index.pdf}{\faIcon{file-pdf}}
  \href{https://chaska-x.netlify.app/operating-system/2023-06-19-03-instalacion-linux}{03
  Instalacion Linux}
\item
  \href{https://chaska-x.netlify.app/operating-system/2023-06-20-04-administracion-particiones-volumenes/index.pdf}{\faIcon{file-pdf}}
  \href{https://chaska-x.netlify.app/operating-system/2023-06-20-04-administracion-particiones-volumenes}{04
  Administracion Particiones Volumenes}
\item
  \href{https://chaska-x.netlify.app/operating-system/2023-07-01-atajos-de-teclado-y-comandos-para-usar-vim/index.pdf}{\faIcon{file-pdf}}
  \href{https://chaska-x.netlify.app/operating-system/2023-07-01-atajos-de-teclado-y-comandos-para-usar-vim}{Atajos
  De Teclado Y Comandos Para Usar Vim}
\item
  \href{https://chaska-x.netlify.app/operating-system/2024-07-15-instalando-specitify/index.pdf}{\faIcon{file-pdf}}
  \href{https://chaska-x.netlify.app/operating-system/2024-07-15-instalando-specitify}{Instalando
  Specitify}
\item
  \href{https://chaska-x.netlify.app/operating-system/2025-07-10-gestiona-tus-dotfiles-con-gnu-stow/index.pdf}{\faIcon{file-pdf}}
  \href{https://chaska-x.netlify.app/operating-system/2025-07-10-gestiona-tus-dotfiles-con-gnu-stow}{Gestiona
  Tus Dotfiles Con Gnu Stow}
\end{enumerate}

Esperamos que encuentres estas publicaciones igualmente interesantes y
útiles. ¡Disfruta de la lectura!






\end{document}
